% !TeX root=../main.tex
\chapter{پیاده‌سازی عملی، ساخت نمونه و آزمون‌های واقعی}

\section{مقدمه}
\label{sec:chapter5_intro}

در فصل چهارم، چارچوب پیشنهادی ناوبری خودران اسکوتر در محیط شبیه‌سازی طراحی و اجزای اصلی آن، شامل مدل سامانه، لایه ادراک، فضای مشاهده و عمل، معماری عامل یادگیری تقویتی، تابع پاداش، سپر ایمنی و فرایند آموزش تشریح شد. فصل حاضر به انتقال این چارچوب از محیط شبیه‌سازی به نمونه واقعی اسکوتر و پیاده‌سازی یکپارچه اجزای سخت‌افزاری و نرم‌افزاری اختصاص دارد. هدف اصلی، ایجاد سامانه‌ای عملی است که بتواند داده‌های حسگرهای واقعی را دریافت و پردازش کرده، ورودی سازگار با سیاست آموزش‌دیده را تشکیل دهد و فرمان‌های تولیدشده را با رعایت محدودیت‌های حرکتی و ایمنی به عملگرهای اسکوتر اعمال کند.

پیاده‌سازی عملی سامانه مستلزم برقراری ارتباط منسجم میان چند لایه محاسباتی و کنترلی است. در این ساختار، پردازش‌های ادراکی و اجرای مدل روی رایانه و رزبری‌پای انجام می‌شوند و کنترل پایین‌سطح موتورها، خواندن انکودرها و دریافت داده برخی حسگرها بر عهده میکروکنترلرها قرار دارد. انتقال داده و فرمان میان این اجزا از طریق ارتباطات سریال و بستر \lr{ROS2} انجام می‌شود. همچنین، تفاوت میان رفتار ایده‌آل شبیه‌سازی و محدودیت‌های نمونه واقعی، از جمله تأخیر عملگرها، رانش لغزشی، نویز حسگرها و نامعتبرشدن موقت داده‌ها، در طراحی زنجیره اجرایی مورد توجه قرار گرفته است.

در ادامه فصل، ابتدا معماری سامانه عملی و نحوه ارتباط اجزای سخت‌افزاری و نرم‌افزاری معرفی می‌شود. سپس ساخت و تجهیز نمونه واقعی اسکوتر، جانمایی حسگرها، سامانه تغذیه و ساختار ارتباطی اجزا تشریح می‌گردد. پس از آن، روش پردازش تصاویر استریو، تشخیص هدف و موانع، تخمین فاصله و سرعت نسبی، پردازش داده حسگرهای فراصوت و تشکیل مشاهده واقعی ارائه می‌شود. نحوه محاسبه اودومتری نسبی با ترکیب داده‌های انکودر و واحد اندازه‌گیری اینرسی نیز به‌منظور ثبت حرکت و تحلیل مسیر توضیح داده خواهد شد.

در بخش‌های بعد، فرایند بارگذاری و اجرای مدل آموزش‌دیده با وزن‌های ثابت، تبدیل خروجی سیاست به فرمان حرکتی، هموارسازی و محدودسازی فرمان، اعمال سپر ایمنی و کنترل سرعت چهار موتور بررسی می‌شود. در پایان، مراحل آماده‌سازی سامانه، سناریوهای آزمون عملی و داده‌های ثبت‌شده برای ارزیابی عملکرد معرفی می‌گردند. نتایج حاصل از آزمون اجزای منفرد و عملکرد یکپارچه سامانه در فصل ششم تحلیل خواهند شد.


\section{معماری سامانه عملی اسکوتر خودران}
\label{sec:real_scooter_architecture}

\subsection{نمای کلی معماری سه‌سطحی سامانه}

این فصل روند پیاده‌سازی سامانه پیشنهادی، آماده‌سازی نمونه واقعی اسکوتر، ارتباط واحدهای سخت‌افزاری و نرم‌افزاری و اجرای آزمون‌های عملی را تشریح می‌کند. هدف، انتقال ساختار طراحی‌شده در شبیه‌سازی به سامانه‌ای فیزیکی و قابل آزمون است؛ به‌گونه‌ای که داده حسگرهای واقعی به مشاهده عامل یادگیری تقویتی و خروجی مدل آموزش‌دیده به فرمان قابل اعمال به موتورهای اسکوتر تبدیل شود.

معماری سامانه عملی سه سطح دارد. در سطح پایین، آردوینو مگا مستقیماً با انکودرهای چهار موتور، راه‌اندازهای موتور و پنج حسگر فراصوت ارتباط دارد، داده چرخ‌ها و فاصله‌های نزدیک‌برد را دریافت می‌کند و فرمان نهایی را به راه‌اندازها می‌فرستد. در سطح میانی، رزبری پای ۵ داده‌های دوربین استریو، واحد اندازه‌گیری اینرسی و آردوینو را جمع‌آوری و برای پردازش به لپ‌تاپ منتقل می‌کند. در سطح بالا، لپ‌تاپ پردازش لحظه‌ای، تشکیل مشاهده، اجرای مدل آموزش‌دیده و تولید فرمان حرکتی را انجام می‌دهد. مدل یادگیری تقویتی در سامانه واقعی فقط در حالت استنتاج اجرا می‌شود و آموزش جدیدی صورت نمی‌گیرد.

این تفکیک، کنترل پایین‌سطح، جمع‌آوری و انتقال داده حسگری، و پردازش و تصمیم‌گیری سطح‌بالا را از یکدیگر جدا می‌کند. خواندن انکودر و حسگرهای فراصوت و اعمال فرمان موتور در سطح پایین باقی می‌ماند، درحالی‌که تخمین عمق، تشخیص هدف و مانع، تشکیل مشاهده و اجرای سیاست روی لپ‌تاپ انجام می‌شود. در نتیجه، بار واحدهای تعبیه‌شده کاهش می‌یابد و هر بخش را می‌توان مستقل توسعه و آزمون کرد.

شکل \ref{fig:practical_system_data_command_flow} جریان کلی داده و فرمان را نشان می‌دهد؛ مسیر داده با فلش آبی و فرمان با فلش نارنجی مشخص شده است. تصاویر استریو و داده واحد اندازه‌گیری اینرسی وارد رزبری پای می‌شوند، اما داده انکودر و حسگرهای فراصوت ابتدا در آردوینو دریافت و از طریق \lr{USB} به رزبری پای منتقل می‌گردند. داده‌های تجمیع‌شده از شبکه به لپ‌تاپ می‌رسند و فرمان تولیدشده نیز در مسیر معکوس، از لپ‌تاپ به رزبری پای، آردوینو، راه‌اندازها و موتورهای اسکوتر اعمال می‌شود.

\begin{figure}[ht]
		\centerline{\includegraphics[width=0.75\linewidth]{robot_sensor_and_control_system_diagram.png}}
	\caption{جریان کلی داده و فرمان در سامانه عملی اسکوتر خودران}
	\label{fig:practical_system_data_command_flow}
	\vspace{0.4em}
	
	{\small
		\href{https://mohammadabedi1179.github.io/thesis-archive/chapters/ch5/sensor-control-system-diagram.html}
		{\textbf{مشاهده نسخه تعاملی شکل با کیفیت بالاتر}}
		\qquad
		\qrcode[height=1.2cm]{https://mohammadabedi1179.github.io/thesis-archive/chapters/ch5/sensor-control-system-diagram.html}
	}
\end{figure}

\subsection{ساختار پردازشی و کنترلی سامانه}
\label{subsec:processing_control_structure}

مسیرهای ارتباطی سامانه به انتقال تصویر و انتقال داده‌های عددی و فرمان تقسیم شدند. تصاویر حجیم استریو به‌صورت دو جریان چپ و راست با \lr{GStreamer} روی شبکه محلی ارسال می‌شوند، درحالی‌که اودومتری، مقادیر فراصوت، سرعت موتورها و فرمان‌های حرکتی با بسته‌های \lr{UDP} میان رزبری پای و لپ‌تاپ انتقال می‌یابند.

رزبری پای افزون بر انتقال داده، بخشی از پردازش حرکتی را انجام می‌دهد: داده انکودر و فراصوت را از آردوینو و \lr{yaw} را از واحد اندازه‌گیری اینرسی دریافت می‌کند، اودومتری نسبی را می‌سازد و آن را همراه با فاصله‌های فراصوت و سرعت موتورها برای لپ‌تاپ می‌فرستد. در جهت مقابل، لپ‌تاپ فرمان مدل را به سرعت هدف چهار موتور تبدیل و با \lr{UDP} ارسال می‌کند؛ رزبری پای نیز آن را به قالب سریال آردوینو تبدیل کرده و از \lr{USB/Serial} می‌فرستد.

از دید اجرایی، مدل \lr{LSTM-SAC} روی لپ‌تاپ سرعت خطی و زاویه‌ای را تولید می‌کند؛ گره واسط سخت‌افزاری آن را به \lr{RPM} چهار موتور تبدیل و با \lr{UDP} به رزبری پای ارسال می‌کند. فرمان سپس در قالب \serialcmd به آردوینو می‌رسد تا کنترل پایین‌سطح سرعت چهار موتور را با بازخورد انکودر انجام دهد.

جدول \ref{tab:real_system_data_command_paths} مسیرهای اصلی تبادل داده و فرمان را خلاصه می‌کند.

\begin{table}[htbp]
	\centering
	\caption{مسیرهای اصلی تبادل داده و فرمان در سامانه عملی}
	\label{tab:real_system_data_command_paths}
	\begin{tabular}{|p{0.2\textwidth}|p{0.58\textwidth}|}
		\hline
		\textbf{مسیر ارتباطی} &
		\textbf{محتوای منتقل‌شده و بستر ارتباطی}\\
		\hline
		
		دوربین استریو به لپ‌تاپ &
		انتقال جریان‌های تصویری چپ و راست با \lr{GStreamer} روی \lr{LAN}. \\
		\hline
		
		آردوینو به رزبری پای &
		ارسال داده انکودرهای چهار موتور و مقادیر حسگرهای فراصوت از طریق \lr{USB/Serial}. \\
		\hline
		
		رزبری پای به لپ‌تاپ &
		ارسال تله‌متری شامل اودومتری نسبی، داده‌های فراصوت و سرعت موتورها با \lr{UDP} روی \lr{LAN}. \\
		\hline
		
		لپ‌تاپ به رزبری پای &
		ارسال فرمان سرعت هدف چهار موتور پس از تبدیل \lr{cmd\_vel} به \lr{RPM}، با \lr{UDP} روی \lr{LAN}. \\
		\hline
		
		رزبری پای به آردوینو &
		ارسال فرمان نهایی چهارموتوره در قالب \serialcmd{} از طریق \lr{USB/Serial}. \\
		\hline
	\end{tabular}
	
\end{table}

\subsection{لایه حسگری و انتقال داده‌ها}
\label{subsec:sensing_and_data_transfer}

لایه حسگری اطلاعات لازم برای ناوبری هدف‌محور، تشخیص موانع، پایش نزدیک‌برد و تخمین حرکت اسکوتر را فراهم می‌کند. منابع اصلی شامل دوربین استریو، پنج حسگر فراصوت، انکودرهای موتور و واحد اندازه‌گیری اینرسی هستند و خروجی آن‌ها در تشکیل مشاهده، اودومتری، ایمنی و ثبت آزمون به کار می‌رود.

دوربین استریو از دو حسگر \lr{IMX219-83} متصل به رزبری پای از طریق \lr{CSI} تشکیل شده است. تصاویر چپ و راست با وضوح ($1640 \times 1232$)، پس از استخراج صفحه روشنایی، به‌صورت خاکستری \lr{mono8} و با نرخ تقریبی ۱۵ فریم بر ثانیه از طریق \lr{GStreamer} و \lr{LAN} به لپ‌تاپ می‌رسند. این تصاویر برای تصحیح هندسی، اختلاف مکانی، عمق، تشخیص هدف و مانع استفاده می‌شوند.

پنج حسگر \lr{SRF05} در موقعیت‌های جلو-چپ، جلو-راست، چپ، راست و عقب نصب و به آردوینو متصل شدند. فاصله‌ها از مسیر \lr{USB/Serial} به رزبری پای و سپس با تله‌متری \lr{UDP} به لپ‌تاپ می‌رسند و هم در مشاهده عامل و هم در منطق ایمنی نزدیک‌برد استفاده می‌شوند.

هر موتور انکودر مستقلی دارد. آردوینو با شمارش پالس‌ها سرعت چهار چرخ را محاسبه می‌کند و این داده را برای کنترل پایین‌سطح و تخمین سرعت حرکتی اسکوتر در اختیار رزبری پای قرار می‌دهد.

به‌دلیل ساختار \lr{Skid-steering}، انکودرها به‌تنهایی در چرخش و لغزش جهت‌گیری را دقیق تخمین نمی‌زنند. ازاین‌رو حسگر \lr{TSS-USB-S v2.0} از طریق \lr{USB} به رزبری پای متصل شد تا زاویه \lr{yaw} را فراهم کند؛ مقدار ابتدای هر آزمون نیز مرجع صفر در نظر گرفته می‌شود.

رزبری پای با ترکیب انکودر و \lr{yaw}، اودومتری نسبی شامل موقعیت، جهت‌گیری و سرعت‌های خطی و زاویه‌ای را می‌سازد و از طریق \lr{UDP} به لپ‌تاپ می‌فرستد. بنابراین داده خام \lr{IMU} مستقیماً وارد مشاهده نمی‌شود؛ خروجی ترکیبی برای تشکیل مشاهده و مدیریت حافظه موانع استفاده می‌گردد.

برای جلوگیری از تداخل، جریان‌های تصویری پرحجم با \lr{GStreamer} و داده‌های عددی مانند اودومتری، فراصوت و سرعت موتورها با بسته‌های \lr{UDP} روی شبکه محلی منتقل می‌شوند.

جدول \ref{tab:real_sensor_data_paths} منابع حسگری و مسیر انتقال آن‌ها را خلاصه می‌کند.

\begin{table}[htbp]
	\centering
	\caption{خلاصه منابع حسگری و مسیر انتقال داده در سامانه عملی}
	\label{tab:real_sensor_data_paths}
	\begin{tabular}{|p{0.24\textwidth}|p{0.31\textwidth}|p{0.35\textwidth}|}
		\hline
		\textbf{منبع داده} &
		\textbf{داده اصلی} &
		\textbf{مسیر انتقال و کاربرد} \\
		\hline
		دوربین استریو &
		تصاویر خاکستری چپ و راست &
		اتصال به رزبری پای از طریق \lr{CSI} و انتقال به لپ‌تاپ با \lr{GStreamer} برای تخمین عمق، تشخیص هدف و تشخیص موانع. \\
		\hline
		
		حسگرهای فراصوت &
		فاصله نزدیک‌برد در پنج جهت &
		اتصال به آردوینو و انتقال از مسیر رزبری پای به لپ‌تاپ برای مشاهده عامل و لایه ایمنی. \\
		\hline
		
		انکودرهای موتور &
		سرعت چرخ‌های چهارگانه &
		خواندن توسط آردوینو و استفاده در کنترل سرعت موتور و تخمین حرکت اسکوتر. \\
		\hline
		
		واحد اندازه‌گیری اینرسی &
		زاویه نسبی \lr{yaw} &
		اتصال به رزبری پای از طریق \lr{USB} و ترکیب با داده انکودر برای تولید اودومتری نسبی. \\
		\hline
		
	\end{tabular}
	
	
\end{table}

\subsection{جریان فرمان در سامانه عملی}
\label{subsec:real_command_flow}

جریان فرمان از مدل \lr{LSTM-SAC} روی لپ‌تاپ آغاز می‌شود. مدل با دریافت مشاهده، کنش دوبعدی نرمال‌شده‌ای تولید می‌کند که طبق رابطه \ref{eq:action_scaling_final} به سرعت خطی و زاویه‌ای مطلوب تبدیل می‌شود؛ بنابراین خروجی مستقیم مدل نه \lr{PWM} و نه فرمان مستقل موتورها، بلکه فرمان سطح‌بالای $(v,\omega)$ است.

فرمان $(v,\omega)$ در لپ‌تاپ به سرعت هدف چهار موتور تبدیل و با \lr{UDP} به رزبری پای ارسال می‌شود. رزبری پای آن را به قالب سریال قابل فهم آردوینو تبدیل و از طریق \lr{USB/Serial} منتقل می‌کند.

آردوینو مقادیر هدف را استخراج و حلقه کنترل سرعت هر موتور را با بازخورد انکودر می‌بندد؛ سپس فرمان مناسب راه‌اندازها را تولید می‌کند. بنابراین نقش آن صرفاً انتقال فرمان نیست.

محدودیت‌های ایمنی نیز در این مسیر اعمال می‌شوند. در حالت عادی فرمان عامل اجرا می‌شود، اما در وضعیت خطر می‌تواند محدود یا متوقف گردد. آستانه‌ها مطابق بخش \ref{subsec:safety_shield_levels_final} و جدول \ref{tab:shield_thresholds_final} هستند و اجرای عملی در بخش \ref{subsec:real_safety_logic_execution} تشریح می‌شود.

برای جلوگیری از ماندن آخرین فرمان هنگام خطا، سازوکار توقف خودکار در لپ‌تاپ و رزبری پای در نظر گرفته شد. اگر فرمان تازه‌ای از مدل یا لپ‌تاپ دریافت نشود، فرمان صفر به موتورها ارسال می‌شود تا قطع ارتباط، توقف نرم‌افزار یا تأخیر غیرعادی موجب ادامه حرکت نشود.

\subsection{ارتباط معماری واقعی با معماری شبیه‌سازی}
\label{subsec:sim_real_architecture_relation}

معماری سامانه عملی بر پایه همان منطق کلی حلقه ادراک، تصمیم‌گیری و کنترل طراحی شد که در فصل چهارم برای محیط شبیه‌سازی معرفی شده بود. در هر دو بستر، ابتدا داده‌های محیطی دریافت و پردازش می‌شوند، سپس مشاهده عامل ساخته می‌شود، عامل تصمیم‌گیر فرمان حرکتی را تولید می‌کند و در نهایت فرمان پس از اعمال محدودیت‌های حرکتی و ایمنی به سامانه حرکتی منتقل می‌شود. حفظ این ساختار برای انتقال مدل آموزش‌دیده از شبیه‌سازی به سامانه واقعی ضروری است، زیرا عامل یادگیری تقویتی با قالب مشخصی از مشاهده و عمل آموزش دیده و در اجرای عملی نیز باید داده‌هایی با مفهوم، ترتیب، مقیاس و محدوده سازگار دریافت کند.

هدف، بازسازی همان زنجیره در سخت‌افزار واقعی بود، نه تغییر منطق اصلی آن. تفاوت عمده در منبع و کیفیت داده‌هاست: حسگرهای واقعی در مقایسه با مدل‌های شبیه‌سازی دارای نویز، تأخیر، محدودیت میدان دید، خطای کالیبراسیون و محدودیت نرخ پردازش‌اند.

یکی از تفاوت‌های مهم، روش تشخیص هدف است. در محیط شبیه‌سازی، هدف بر اساس ویژگی رنگی آن از تصویر جدا می‌شد؛ اما در سامانه عملی، تصاویر استریو به‌صورت خاکستری منتقل می‌شوند و اطلاعات رنگی برای تشخیص هدف در دسترس نیست. به همین دلیل، در اجرای واقعی از نشانگر \lr{AprilTag} استفاده شد. پس از تشخیص نشانگر، فاصله و زاویه نسبی هدف از ترکیب موقعیت تگ در تصویر و نقشه اختلاف مکانی استخراج می‌شود. سپس مشابه روشی که برای موانع استفاده می‌شود، مختصات نسبی هدف در زمان‌های متوالی با فیلتر کالمن پالایش شده و سرعت نسبی هدف محاسبه می‌گردد. بنابراین، اگرچه روش استخراج هدف از تصویر در سامانه عملی نسبت به شبیه‌سازی تغییر کرده است، خروجی نهایی همچنان شامل فاصله، زاویه و مؤلفه‌های سرعت نسبی هدف در قالب مورد نیاز عامل است.

تفاوت مهم دیگر به تخمین وضعیت حرکتی اسکوتر مربوط است. در محیط شبیه‌سازی، موقعیت و جهت‌گیری نسبی ربات از خروجی اودومتری و وضعیت منتشرشده توسط \lr{Gazebo} قابل دریافت بود. اما در سامانه واقعی، چنین مرجع مستقیمی وجود ندارد و موقعیت نسبی اسکوتر باید از داده‌های حسگرهای نصب‌شده روی سامانه تخمین زده شود. به همین دلیل، در پیاده‌سازی، سرعت حرکتی از داده انکودرهای چرخ‌ها و جهت‌گیری از زاویه \lr{yaw} واحد اندازه‌گیری اینرسی استخراج شده و با ترکیب این داده‌ها، اودومتری نسبی اسکوتر نسبت به نقطه شروع آزمون محاسبه می‌شود. بنابراین، در حالی که در شبیه‌سازی موقعیت نسبی از وضعیت شبیه‌ساز به دست می‌آید، در سامانه واقعی این کمیت حاصل یک تخمین حسگری است و می‌تواند تحت تأثیر لغزش چرخ‌ها، نویز \lr{IMU} و خطای تجمعی انتگرال‌گیری قرار گیرد.

همچنین تفاوت میان مدل حرکتی شبیه‌سازی و سامانه فیزیکی واقعی باید در نظر گرفته شود. در شبیه‌سازی، پاسخ موتورها و اعمال فرمان به چرخ‌ها به‌صورت ایده‌آل‌تر و با محدودیت‌های کمتر مدل می‌شود، در حالی که در سامانه واقعی، ویژگی‌هایی مانند دینامیک موتور، تأخیر راه‌اندازها، محدودیت باتری، اشباع \lr{PWM}، اصطکاک، لقی مکانیکی و لغزش چرخ‌ها بر اجرای فرمان اثر می‌گذارند. به همین دلیل، در سامانه عملی فرمان خروجی عامل مستقیماً به موتورها اعمال نمی‌شود، بلکه ابتدا به سرعت هدف چهار موتور تبدیل شده و سپس آردوینو با استفاده از بازخورد انکودرها، کنترل پایین‌سطح سرعت موتورها را انجام می‌دهد. این لایه موجب می‌شود اثر بخشی از تفاوت‌های میان مدل شبیه‌سازی و سخت‌افزار واقعی کاهش یابد.

برای موانع نیز ساختار اصلی حفظ شد: مانع در تصویر تشخیص داده می‌شود، فاصله از نقشه اختلاف مکانی به دست می‌آید و موقعیت و سرعت نسبی با فیلتر کالمن پالایش می‌گردد. بدین‌ترتیب داده خام هدف و موانع پیش از ورود به مشاهده، پایدار و با ساختار آموزش سازگار می‌شود.

جدول \ref{tab:sim_real_mapping_updated} نگاشت اجزای اصلی معماری شبیه‌سازی به سامانه عملی را نشان می‌دهد.

\begin{table}[htb]
	\centering
	\caption{نگاشت خلاصه اجزای معماری شبیه‌سازی به سامانه عملی}
	\label{tab:sim_real_mapping_updated}
	\renewcommand{\arraystretch}{1.18}
	\begin{tabular}{|p{0.23\textwidth}|p{0.31\textwidth}|p{0.36\textwidth}|}
		\hline
		\textbf{بخش سامانه} &
		\textbf{در محیط شبیه‌سازی} &
		\textbf{در سامانه عملی} \\
		\hline
	
		ادراک دیداری و عمق &
		تصاویر استریو، تصحیح هندسی و تولید نقشه اختلاف مکانی با پارامترهای شبیه‌سازی &
		تصاویر دوربین واقعی، کالیبراسیون عملی، تصحیح هندسی و تولید نقشه اختلاف مکانی با همان منطق پردازشی \\
		\hline
		
		تشخیص و پالایش هدف &
		جداسازی هدف بر اساس ویژگی رنگی و استخراج فاصله، زاویه و سرعت نسبی هدف &
		تشخیص هدف با \lr{AprilTag}، استخراج فاصله و زاویه نسبی، و محاسبه سرعت نسبی هدف با فیلتر کالمن \\
		\hline
		
		تشخیص و پالایش موانع &
		تشخیص موانع، تخمین فاصله و پالایش زمانی داده‌ها در محیط شبیه‌سازی &
		تشخیص موانع از تصویر واقعی، تخمین عمق از نقشه اختلاف مکانی و پالایش موقعیت و سرعت نسبی موانع با فیلتر کالمن \\
		\hline
		
		وضعیت حرکتی و اودومتری &
		دریافت موقعیت، جهت‌گیری و سرعت از خروجی اودومتری و وضعیت \lr{Gazebo} &
		تخمین اودومتری نسبی با ترکیب سرعت حاصل از انکودرهای چرخ و زاویه \lr{yaw} واحد اندازه‌گیری اینرسی \\
		\hline
		
		اعمال فرمان به موتورها &
		اعمال فرمان به مدل حرکتی شبیه‌سازی‌شده با دینامیک ساده‌تر و محدودیت‌های کمتر &
		تبدیل فرمان عامل به سرعت هدف چهار موتور و اجرای آن با کنترل پایین‌سطح مبتنی بر بازخورد انکودر \\
		\hline
		
		حسگرهای نزدیک‌برد &
		مدل حسگرهای فراصوت برای مشاهده و سپر ایمنی &
		پنج حسگر فراصوت واقعی برای مشاهده عامل و تشخیص وضعیت‌های نزدیک‌برد \\
		\hline
		
		تشکیل مشاهده و اجرای مدل &
		ساخت مشاهده از داده‌های شبیه‌سازی و اجرای مدل آموزش‌دیده &
		تبدیل داده‌های واقعی به همان قالب مشاهده و اجرای همان مدل \lr{LSTM-SAC} در حالت استنتاج \\
		\hline
		
		فرمان و ایمنی &
		تولید فرمان حرکتی و اعمال منطق سپر ایمنی در محیط شبیه‌سازی &
		تبدیل فرمان عامل به فرمان قابل اجرا روی سخت‌افزار و اعمال محدودیت‌های ایمنی در مسیر اجرا \\
		\hline
		
	\end{tabular}

	
\end{table}

سامانه عملی، پیاده‌سازی فیزیکی معماری شبیه‌سازی است، اما انتقال مدل به‌تنهایی کافی نیست. واسط‌های ادراکی، ارتباطی و کنترلی باید داده خام حسگر را به مشاهده سازگار با عامل و خروجی آن را به فرمان موتور تبدیل کنند. بنابراین منطق ادراک، مشاهده، تصمیم‌گیری، کنترل و ایمنی حفظ شده و جزئیات اجرا متناسب با نویز، محدودیت و نرخ پردازش سخت‌افزار تنظیم شده‌اند.

\section{تجهیز و آماده‌سازی سخت‌افزاری سامانه عملی}
\label{sec:real_hardware_preparation}

پس از تعیین معماری، نمونه فیزیکی برای نصب حسگرها، واحدهای پردازشی و اجزای کنترلی آماده شد. این بخش جانمایی و نصب سخت‌افزار را بیان می‌کند و پردازش داده، تشکیل مشاهده و اجرای مدل در بخش‌های بعدی می‌آید.

\subsection{بستر مکانیکی اسکوتر و تغییرات اعمال‌شده}
\label{subsec:real_scooter_platform}

بستر مکانیکی، یک اسکوتر چهارچرخ موجود با چهار موتور، راه‌انداز و ساختار حرکت مبتنی بر لغزش بود؛ ازاین‌رو طراحی شاسی، انتخاب موتور و توسعه اولیه سامانه حرکتی در این پژوهش انجام نشد. شاسی و مجموعه حرکتی بدون تغییر اساسی حفظ و آماده‌سازی به نصب حسگرها، واحدهای پردازشی و کنترلی و مسیرهای توان و سیگنال محدود شد.

داخل بدنه برای آردوینو مگا، راه‌اندازها، باتری‌ها و کابل‌های توان و سیگنال استفاده شد. در آزمون آزمایشگاهی، لپ‌تاپ برای دسترسی نرم‌افزاری، ثبت داده و عیب‌یابی روی درپوش فوقانی قرار گرفت؛ در نسخه نهایی می‌توان آن را با واحدی کوچک‌تر و نصب‌شونده جایگزین کرد.

\subsection{نصب دوربین استریو و رزبری پای}
\label{subsec:stereo_camera_raspberry_installation}

دوربین استریو در بخش جلویی اسکوتر نصب شد تا میدان دید آن با جهت حرکت اصلی اسکوتر هم‌راستا باشد. برای نصب دوربین، ابتدا یک پایه فلزی به ارتفاع ۷.۰ متر به قسمت جلویی بدنه پیچ شد. سپس یک صفحه از جنس \lr{ABS} به‌صورت موازی با پایه فلزی روی آن نصب گردید و ماژول دوربین استریو به این صفحه پیچ شد. این ساختار، سطح نصب مناسبی برای دوربین فراهم کرد و موقعیت آن را نسبت به بدنه ثابت نگه داشت.

رزبری پای ۵ نیز در نزدیکی دوربین و روی همان مجموعه نصب جلویی قرار گرفت. برای این منظور، در پشت پایه فلزی، یک صفحه \lr{ABS} دیگر به‌صورت عمودی و با استفاده از براکت به پایه متصل شد و رزبری پای روی این صفحه پیچ گردید. فاصله رزبری پای تا دوربین در حدود ۵ سانتی‌متر است؛ بنابراین اتصال دو دوربین به رزبری از طریق کابل‌های \lr{CSI} بدون کابل‌کشی بلند انجام می‌شود.

\begin{figure}[htbp]
	\centering
	\includegraphics[width=0.75\linewidth]{camera_and_raspberry}
	\caption{تصویر دوربین استریو و رزبری نصب شده بر روی پایه فلزی اسکوتر}
	\label{fig:camera_and_raspberry}
\end{figure}

\subsection{نصب حسگرهای فراصوت}
\label{subsec:ultrasonic_sensor_installation}

پنج حسگر فراصوت \lr{SRF05} در موقعیت‌های جلو-چپ، جلو-راست، چپ، راست و عقب اسکوتر نصب شدند. ارتفاع نصب حسگرها از سطح زمین حدود \(15\) سانتی‌متر در نظر گرفته شد. محل و جهت‌گیری هر حسگر نیز مطابق نواحی ادراکی تعریف‌شده در فصل چهارم انتخاب شد تا داده‌های اندازه‌گیری‌شده در سامانه واقعی با ساختار فضای مشاهده و منطق عملکرد سپر ایمنی سازگار باشند.

به‌منظور تثبیت حسگرها روی بدنه اسکوتر و محافظت از آن‌ها در برابر ضربه و تماس مستقیم، یک نگهدارنده اختصاصی طراحی شد. این نگهدارنده دارای دو خار تثبیت‌کننده است که در حفره‌های موجود روی برد حسگر قرار می‌گیرند و از جابه‌جایی حسگر نسبت به نگهدارنده جلوگیری می‌کنند. همچنین، یک شیلد محافظ برای پوشاندن برد و قطعات الکترونیکی حسگر در نظر گرفته شد. در طراحی این شیلد، دهانه‌های فرستنده و گیرنده فراصوت کاملاً باز نگه داشته شدند تا پوشش محافظ موجب انسداد مسیر انتشار و دریافت امواج یا ایجاد اختلال در اندازه‌گیری فاصله نشود.

در شکل \ref{fig:ultrasonic_sensor_design_and_installation}، مدل حسگر بدون شیلد، مدل مجهز به شیلد محافظ و یک نمونه از حسگر نصب‌شده روی بخش عقب اسکوتر نشان داده شده است. هنگام نصب، زاویه هر حسگر به‌گونه‌ای تنظیم شد که ناحیه موردنظر را پوشش دهد و میدان دید آن توسط بدنه اسکوتر، نگهدارنده یا سایر تجهیزات جانبی مسدود نشود.

\begin{figure}[htbp]
	\centering
	
	\subfloat[]{
		\includegraphics[width=0.47\textwidth]
		{SRF05.png}
	}
	\hfill
	\subfloat[]{
		\includegraphics[width=0.47\textwidth]
		{SRF05_shield.png}
	}
	
	\vspace{0.3cm}
	
	\subfloat[]{
		\includegraphics[width=0.5\textwidth]
		{ultrasoinc_applied}
	}
	
	\caption{طراحی و نصب حسگر فراصوت \lr{SRF05}:
		الف) مدل حسگر بدون شیلد محافظ،
		ب) مدل حسگر مجهز به شیلد محافظ،
		ج) نمونه حسگر نصب‌شده روی بخش عقب اسکوتر}
	\label{fig:ultrasonic_sensor_design_and_installation}
\end{figure}

\subsection{جانمایی واحد اندازه‌گیری اینرسی، آردوینو، راه‌اندازهای موتور و سامانه تغذیه}
\label{subsec:imu_arduino_driver_power_installation}

واحد اندازه‌گیری اینرسی از طریق اتصال \lr{USB} به رزبری پای متصل شد. آردوینو در داخل محفظه اصلی بدنه نصب شد و ارتباط با انکودرهای چهار موتور، حسگرهای فراصوت و راه‌اندازهای موتور را بر عهده گرفت. راه‌اندازهای موتور نیز داخل بدنه و روی یک صفحه فلزی ثابت شدند. این صفحه علاوه بر ایجاد محل نصب مکانیکی، به دفع حرارت تولیدشده توسط راه‌اندازها کمک می‌کند. راه‌اندازها سیگنال کنترلی را از آردوینو دریافت کرده و توان لازم را به موتورهای چرخ منتقل می‌کنند.

سامانه تغذیه از چهار باتری لیتیوم‌یونی \lr{Pleson 10S2P} تشکیل شده است. مشخصات هر باتری شامل ولتاژ نامی $36~V$، ظرفیت $4.4~Ah$، انرژی $158.4~Wh$ و ولتاژ شارژ $42~V$ است. باتری‌ها به‌صورت موازی بسته شده‌اند و باس توان $36~V$ را تشکیل می‌دهند. این باس برای تغذیه راه‌اندازهای موتور استفاده می‌شود و از طریق تنظیم ولتاژ، تغذیه $5~V$ موردنیاز رزبری پای، آردوینو و حسگرهای وابسته نیز فراهم می‌گردد. لپ‌تاپ در نسخه آزمایشگاهی از منبع تغذیه مستقل خود استفاده می‌کند.

شکل \ref{fig:power_and_signal_distribution} مسیرهای توان و سیگنال سامانه را نشان می‌دهد.

\begin{figure}[htbp]

\end{figure}

\begin{figure}[ht]
	\centering
\includegraphics[width=0.75\linewidth]{schematic_of_electronic_power_system}
\caption{شماتیک مسیرهای توان و سیگنال در سامانه عملی اسکوتر خودران}
\label{fig:power_and_signal_distribution}
	\vspace{0.4em}
	{\small
		\href{https://mohammadabedi1179.github.io/thesis-archive/chapters/ch5/schematic_of_electronic_power_system.html}
		{\textbf{مشاهده نسخه تعاملی شکل با کیفیت بالاتر}}
		\qquad
		\qrcode[height=1.2cm]{https://mohammadabedi1179.github.io/thesis-archive/chapters/ch5/schematic_of_electronic_power_system.html}
	}
\end{figure}

پس از جانمایی اجزای مکانیکی، حسگرها، واحدهای پردازشی، باتری‌ها و راه‌اندازها، مدل سه‌بعدی یکپارچه برای بررسی مونتاژ نهایی تهیه شد. شکل \ref{fig:assembled_scooter_overview} شاسی، تعلیق و چرخ‌ها، اجزای داخل بدنه و نصب دوربین استریو و رزبری پای را از چند زاویه نشان می‌دهد؛ نمای بالا نیز توزیع تجهیزات در محفظه اصلی را آشکار می‌کند.

\begin{figure}[ht]
	\centering

\subfloat[]{
	\includegraphics[width=0.47\textwidth]
	{Scooter_3d_3corner_view_1.png}
}
\hfill
\subfloat[]{
	\includegraphics[width=0.47\textwidth]
	{Scooter_3d_3corners_view_2.png}
}

\vspace{0.3cm}

\subfloat[]{
	\includegraphics[width=0.47\textwidth]
	{Scooter_3d_rear_view.png}
}
\hfill
\subfloat[]{
	\includegraphics[width=0.47\textwidth]
	{Scooter_3d_upper_view.png}
}

\caption{نمای کلی مدل سه‌بعدی مونتاژشده اسکوتر و جانمایی اجزای
	سخت‌افزاری: الف) نمای سه‌بعدی از سمت جلو و راست،
	ب) نمای سه‌بعدی از سمت جلو و چپ،
	ج) نمای روبه‌رو از مجموعه تعلیق، چرخ‌ها و پایه تجهیزات،
	د) نمای بالا از جانمایی باتری‌ها، راه‌اندازهای موتور،
	واحدهای کنترلی و تجهیزات پردازشی}
\label{fig:assembled_scooter_overview}
	\vspace{0.4em}
	
	{\small
		\href{https://mohammadabedi1179.github.io/thesis-archive/viewer/scooter-3d-model.html}
		{\textbf{مشاهده نسخه تعاملی مدل سه بعدی اسکوتر}}
		\qquad
		\qrcode[height=1.2cm]{https://mohammadabedi1179.github.io/thesis-archive/viewer/scooter-3d-model.html}
	}
\end{figure}

\section{پیاده‌سازی لایه ادراک در سامانه عملی}
\label{sec:real_perception_layer}

پس از آماده‌سازی سخت‌افزار، لایه ادراک داده خام دوربین استریو، فراصوت، انکودر و واحد اندازه‌گیری اینرسی را به اطلاعات ناوبری تبدیل می‌کند. خروجی آن شامل فاصله و زاویه هدف، موقعیت نسبی موانع، فاصله‌های نزدیک‌برد و وضعیت نسبی حرکت اسکوتر است که برای تشکیل مشاهده عامل استفاده می‌شوند.

در پیاده‌سازی، تلاش شد منطق کلی لایه ادراک با ساختار تعریف‌شده در فصل چهارم سازگار باقی بماند. با این حال، به دلیل استفاده از حسگرهای واقعی، پارامترهای کالیبراسیون، نرخ پردازش، شرایط نوری و نویز اندازه‌گیری با محیط شبیه‌سازی متفاوت هستند. همچنین در تشخیص هدف، به‌جای جداسازی رنگی مورد استفاده در شبیه‌سازی، از نشانگر \lr{AprilTag} استفاده شد؛ زیرا تصاویر استریو در سامانه عملی به‌صورت خاکستری منتقل می‌شوند و اطلاعات رنگی برای تشخیص هدف در دسترس نیست.

\subsection{دریافت تصویر از دوربین استریو}
\label{subsec:real_stereo_image_acquisition}

برای دریافت اطلاعات دیداری محیط، از یک دوربین استریو مبتنی بر دو حسگر \lr{IMX219-83} استفاده شد. این دو حسگر به‌صورت چپ و راست عمل کرده و از طریق دو کابل \lr{CSI} به رزبری متصل هستند. در سامانه عملی، هر دو تصویر ابتدا روی رزبری دریافت می‌شوند و سپس برای پردازش‌های بعدی به لپ‌تاپ منتقل می‌گردند. هدف از این مرحله، فراهم‌سازی یک جفت تصویر هم‌زمان از محیط روبه‌روی اسکوتر است تا در ادامه بتوان از آن‌ها برای تصحیح هندسی، محاسبه اختلاف مکانی، تخمین عمق، تشخیص هدف و تشخیص موانع استفاده کرد.

تصاویر با وضوح ($1640 \times 1232$) و نرخ حدود ۱۵ فریم بر ثانیه از طریق کابل \lr{LAN} و \lr{GStreamer} به لپ‌تاپ انتقال می‌یابند. این ساختار پایداری جریان چپ و راست را با بار ارتباطی و پردازشی قابل قبول فراهم کرد.

در پیاده‌سازی، فریم‌های خام دوربین در قالب \lr{YUV420} دریافت می‌شوند. در این قالب، هر تصویر از سه مؤلفه اصلی تشکیل شده است: صفحه \lr{Y} که شدت روشنایی را نگه می‌دارد، و دو صفحه \lr{U} و \lr{V} که اطلاعات رنگی را بیان می‌کنند. چنان‌که در شکل \ref{fig:real_stereo_image_pipeline} نشان داده شده است، در این پژوهش تنها صفحه \lr{Y} از هر تصویر استخراج و برای ادامه پردازش نگه‌داری می‌شود و صفحات \lr{U} و \lr{V} حذف می‌شوند. دلیل این انتخاب آن است که در پردازش استریو، اطلاعات روشنایی برای تطبیق نقاط متناظر و برآورد اختلاف مکانی کافی است و استفاده از صفحات رنگی نه‌تنها ضرورت مستقیمی ندارد، بلکه باعث افزایش حجم داده و پیچیدگی پردازش نیز می‌شود. در نتیجه، خروجی این مرحله برای هر دوربین یک تصویر خاکستری ۸ بیتی است که ورودی اصلی زنجیره پردازش استریو را تشکیل می‌دهد.

استفاده از صفحه \lr{Y} دو مزیت اصلی دارد. نخست، حجم داده ارسالی و پردازش‌شده نسبت به استفاده از تصویر رنگی کامل کاهش می‌یابد و این موضوع برای اجرای سامانه در زمان واقعی اهمیت دارد. دوم، مرحله تطبیق استریو بر پایه الگوهای شدت روشنایی انجام می‌شود و بنابراین استفاده از تصویر خاکستری برای تولید نقشه اختلاف مکانی کفایت می‌کند. به همین دلیل، در پیاده‌سازی، هر دو تصویر چپ و راست پس از استخراج صفحه \lr{Y}، به‌عنوان ورودی مشترک مراحل بعدی شامل تصحیح هندسی، تطبیق استریو و برآورد عمق استفاده می‌شوند.

بنابراین دریافت تصویر شامل اخذ جفت‌تصویر، تبدیل \lr{YUV420} به روشنایی قابل استفاده و آماده‌سازی ورودی زنجیره ادراک دیداری است، نه صرفاً انتقال داده.

\begin{figure}[htbp]

\end{figure}

\begin{figure}[ht]
	\centering
\includegraphics[width=0.75\linewidth]{stereo_vision_y_plane_extraction_process.png}
\caption{روند دریافت و آماده‌سازی تصاویر دوربین استریو در سامانه عملی؛ از دریافت فریم خام تا استخراج صفحه روشنایی \lr{Y} و استفاده از آن در پردازش استریو}
\label{fig:real_stereo_image_pipeline}
	
	\vspace{0.4em}
	
	{\small
		\href{https://mohammadabedi1179.github.io/thesis-archive/chapters/ch5/stereo_vision_y_plane_extraction_process.html}
		{\textbf{مشاهده نسخه تعاملی شکل با کیفیت بالاتر}}
		\qquad
		\qrcode[height=1.2cm]{https://mohammadabedi1179.github.io/thesis-archive/chapters/ch5/stereo_vision_y_plane_extraction_process.html}
	}
\end{figure}

\subsection{تصحیح هندسی تصاویر و تولید نقشه اختلاف مکانی}
\label{subsec:real_depth_map_generation}

پس از دریافت تصاویر چپ و راست در لپ‌تاپ، ابتدا تصحیح هندسی تصاویر استریو انجام می‌شود. هدف از این مرحله، حذف اثر اعوجاج لنزها و هم‌راستاسازی تصاویر چپ و راست است؛ به‌گونه‌ای که نقاط متناظر در دو تصویر روی خطوط افقی متناظر قرار گیرند. این هم‌راستاسازی برای محاسبه پایدار اختلاف مکانی ضروری است، زیرا در صورت وجود خطای هندسی میان دو تصویر، الگوریتم تطبیق استریو نمی‌تواند جست‌وجوی نقاط متناظر را به‌صورت یک‌بعدی و پایدار انجام دهد.

تصحیح هندسی با استفاده از پارامترهای کالیبراسیون دوربین استریو انجام شد. در کالیبراسیون نهایی، رزولوشن تصویر برابر ($1640 \times 1232$) پیکسل، فاصله پایه دوربین‌ها برابر ($60.172~mm$)، فاصله کانونی تصویر تصحیح‌شده برابر ($1238.075~px$)، و خطای بازفرافکنی استریو برابر ($0.3774~px$) به‌دست آمد. این مقادیر برای ساخت نگاشت‌های تصحیح و همچنین محاسبات فاصله بر اساس اختلاف مکانی استفاده شدند. جزئیات کامل روند کالیبراسیون، ماتریس‌های درونی دوربین‌ها، ضرایب اعوجاج، ماتریس‌های تصحیح و ماتریس‌های تصویر در پیوست \ref{app:stereo_calibration_details} و جدول \ref{tab:stereo_calibration_final_params} ارائه شده‌اند.

در ساختار نرم‌افزاری سامانه، تصحیح هندسی فقط در گره دریافت تصویر انجام می‌شود. این گره پس از دریافت تصاویر چپ و راست، فایل کالیبراسیون را بارگذاری کرده، نگاشت‌های تصحیح را می‌سازد و تصاویر تصحیح‌شده را روی موضوعات ارتباطی مربوط به تصویر چپ و راست منتشر می‌کند. پس از انتشار تصاویر تصحیح‌شده، گره تولید اختلاف مکانی، تصاویر چپ و راست را دریافت کرده و با استفاده از هم‌زمان‌سازی تقریبی، جفت‌تصویرهای نزدیک از نظر زمانی را برای پردازش انتخاب می‌کند. سپس تصاویر خاکستری تصحیح‌شده وارد الگوریتم \lr{StereoSGBM} می‌شوند. این الگوریتم برای هر ناحیه از تصویر چپ، ناحیه متناظر آن را در تصویر راست جست‌وجو کرده و اختلاف افقی میان آن‌ها را به‌عنوان مقدار اختلاف مکانی محاسبه می‌کند. مبنای هندسی تبدیل اختلاف مکانی به عمق همان رابطه ارائه‌شده در فصل چهارم است و در این بخش از تکرار آن خودداری می‌شود؛ برای رابطه عمق بر حسب اختلاف مکانی به رابطه \ref{eq:stereo_depth_basic} مراجعه شود.

در پیاده‌سازی، پارامترهای \lr{SGBM} بر اساس رزولوشن تصویر و محدوده فاصله مورد نیاز در آزمون‌ها تنظیم شدند. مقدار حداقل اختلاف مکانی برابر ۱۶ و مقدار بازه اختلاف مکانی برابر ۱۷۶ در نظر گرفته شد تا بازه مناسبی از اختلاف مکانی برای فاصله‌های کاری سامانه پوشش داده شود. همچنین اندازه بلوک برابر ۷ انتخاب شد تا میان کاهش نویز و حفظ مرزهای عمقی تعادل برقرار شود. از آنجا که محاسبه اختلاف مکانی روی تصویر با وضوح بالا هزینه پردازشی زیادی دارد، در پیاده‌سازی عملی امکان کاهش مقیاس داخلی تصویر پیش از اجرای \lr{SGBM} در نظر گرفته شد. پس از محاسبه اختلاف مکانی در مقیاس کوچک‌تر، مقادیر اختلاف مکانی به واحد پیکسل در وضوح اصلی بازگردانده می‌شوند.

برای بهبود کیفیت خروجی، دو مرحله تکمیلی نیز در زنجیره پردازش قرار گرفت. نخست، از روش \lr{CLAHE} برای بهبود کنتراست موضعی تصاویر خاکستری استفاده شد. این پیش‌پردازش در نواحی کم‌بافت یا دارای تغییرات روشنایی ضعیف می‌تواند به افزایش قابلیت تطبیق میان دو تصویر کمک کند. دوم، از فیلتر \lr{WLS} برای پالایش نقشه اختلاف مکانی استفاده شد. نقشه اختلاف مکانی خام معمولاً شامل نویز، لکه‌های پراکنده و ناپایداری در مرز اجسام است؛ فیلتر \lr{WLS} با استفاده از تصویر مرجع، خروجی تطبیق را هموارتر کرده و در عین حال تلاش می‌کند مرزهای مهم تصویر حفظ شوند. جزئیات کامل الگوریتم تولید نقشه اختلاف مکانی، تنظیمات \lr{SGBM}، نقش \lr{WLS}، اثر \lr{CLAHE} و نحوه تبدیل اختلاف مکانی به عمق در پیوست \ref{appendix:stereo_depth_sgbm} ارائه شده است.

خروجی این مرحله، نقشه اختلاف مکانی متراکم در قالب پیام \lr{DisparityImage} است. در این پیام، علاوه بر تصویر اختلاف مکانی، پارامترهای مورد نیاز برای تفسیر مقادیر اختلاف مکانی، از جمله فاصله کانونی و فاصله پایه دوربین، نیز قرار می‌گیرند. این نقشه مستقیماً وارد فضای مشاهده عامل نمی‌شود، بلکه در مراحل بعدی برای استخراج فاصله هدف و موانع استفاده می‌شود.

شکل \ref{fig:real_disparity_generation_process} نمونه‌ای از مراحل اصلی تولید و پالایش نقشه اختلاف مکانی در سامانه عملی را نشان می‌دهد. در این شکل، ابتدا تصویر ورودی خام نمایش داده شده است. سپس نقشه اختلاف مکانی خام حاصل از اجرای \lr{SGBM}، نقشه پالایش‌شده با \lr{WLS} و خروجی حاصل از پیش‌پردازش \lr{CLAHE} ارائه شده‌اند. این مقایسه نشان می‌دهد که کیفیت نقشه اختلاف مکانی در سامانه واقعی به کیفیت تصویر ورودی، کنتراست موضعی و پالایش پس از تطبیق وابسته است.

\begin{figure}[htbp]
	\centering
	
	\subfloat[]{
		\includegraphics[width=0.48\linewidth]{real_stereo_raw_input.png}
	}
	\hfill
	\subfloat[]{
		\includegraphics[width=0.48\linewidth]{real_disparity_raw.png}
	}
	
	\vspace{0.4cm}
	
	\subfloat[]{
		\includegraphics[width=0.48\linewidth]{real_disparity_wls.png}
	}
	\hfill
	\subfloat[]{
		\includegraphics[width=0.48\linewidth]{real_disparity_clahe.png}
	}
	
	\caption{نمونه مراحل تولید و پالایش نقشه اختلاف مکانی در سامانه عملی:
		الف) تصویر ورودی خام،
		ب) نقشه اختلاف مکانی خام،
		ج) نقشه اختلاف مکانی پس از اعمال \lr{WLS}،
		د) نقشه اختلاف مکانی با پیش‌پردازش \lr{CLAHE}}
	\label{fig:real_disparity_generation_process}
\end{figure}

\subsection{تشخیص هدف با استفاده از نشانگر \lr{AprilTag}}
\label{subsec:real_apriltag_goal_detection}

در محیط شبیه‌سازی، هدف بر اساس ویژگی رنگی آن از تصویر جدا می‌شد. اما در سامانه عملی، چنان‌که در بخش‌های پیشین توضیح داده شد، تصاویر استریو به‌صورت خاکستری منتقل می‌شوند و اطلاعات رنگی برای تشخیص هدف در دسترس نیست. به همین دلیل، در پیاده‌سازی واقعی از نشانگرهای \lr{AprilTag} برای تشخیص هدف استفاده شد. این روش نسبت به جداسازی رنگی وابستگی کمتری به رنگ هدف دارد و در تصویر خاکستری نیز قابل استفاده است.

برای هدف آزمایشگاهی، پنج نشانگر \lr{AprilTag} که در شکل \ref{fig:April_Tags} قابل مشاهده هستند، روی سطوح مختلف جعبه هدف نصب شد. این نشانگرها روی سطوح جانبی و سطح بالایی جعبه قرار گرفتند و روی سطح زیرین جعبه نشانگری نصب نشد. دو سطح حامل تگ دارای ابعاد تقریبی ($50 \times 60~cm$)، دو سطح دیگر دارای ابعاد تقریبی ($30 \times 50~cm$)، و سطح بالایی دارای ابعاد تقریبی ($30 \times 60~cm$) بود. استفاده از چند تگ روی وجوه مختلف هدف موجب شد در زوایا و وضعیت‌های مختلف، حداقل یکی از نشانگرها در میدان دید دوربین قرار گیرد و امکان تشخیص هدف افزایش یابد.

\begin{figure}[ht]
	\centering
\includegraphics[width=0.75\linewidth]{April_Tag_ID_0-->4}
\caption{نشانگرهای استفاده‌شده روی هدف}
\label{fig:April_Tags}
	
	\vspace{0.4em}
	
	{\small
		\href{https://mohammadabedi1179.github.io/thesis-archive/chapters/ch5/April_Tag_ID_0--%3E4.html}
		{\textbf{مشاهده نسخه ۳ بعدی}}
		\qquad
		\qrcode[height=1.2cm]{https://mohammadabedi1179.github.io/thesis-archive/chapters/ch5/April_Tag_ID_0--%3E4.html}
	}
\end{figure}

خانواده نشانگرهای مورد استفاده از نوع \lr{tag36h11} است و شناسه‌های مجاز برای تشخیص هدف در بازه تگ‌های نصب‌شده تعریف شده‌اند. در صورتی که چند نشانگر هم‌زمان در تصویر دیده شوند، نشانگری که سطح تصویر بزرگ‌تری دارد به‌عنوان نماینده هدف انتخاب می‌شود. این انتخاب از آن جهت منطقی است که معمولاً تگ بزرگ‌تر در تصویر، واضح‌تر و پایدارتر قابل تشخیص است و مقدار اختلاف مکانی استخراج‌شده از ناحیه آن نیز قابلیت اطمینان بیشتری دارد.

پس از تشخیص تگ، چهار گوشه آن در تصویر استخراج شده و مرکز تگ محاسبه می‌شود. مرکز تگ برای تعیین زاویه نسبی هدف نسبت به محور دید دوربین به کار می‌رود. یعنی اختلاف افقی مرکز تگ نسبت به مرکز تصویر، همراه با فاصله کانونی دوربین، برای محاسبه زاویه دید هدف استفاده می‌شود. خروجی اولیه این مرحله شامل مختصات مرکز تگ در تصویر، زاویه نسبی هدف و شناسه تگ تشخیص‌داده‌شده است.

برای محاسبه فاصله هدف، پس از تعیین کادر محدودکننده تگ در تصویر چپ، ناحیه متناظر آن روی نقشه اختلاف مکانی انتخاب می‌شود. سپس مقادیر معتبر اختلاف مکانی در این ناحیه استخراج شده و با استفاده از پارامترهای فاصله کانونی و فاصله پایه دوربین، عمق هدف محاسبه می‌شود. مبنای محاسبه عمق همان رابطه اختلاف مکانی ارائه‌شده در فصل چهارم و بخش \ref{subsec:real_depth_map_generation} است.

برای افزایش پایداری تخمین فاصله، کادر تگ به‌صورت محدود گسترش داده می‌شود تا تعداد پیکسل‌های معتبر بیشتری از نقشه اختلاف مکانی در محاسبه استفاده شوند. سپس مقادیر نامعتبر یا کمتر از حد مجاز حذف می‌شوند و در صورتی که تعداد پیکسل‌های معتبر کافی باشد، مقدار نماینده اختلاف مکانی از روی آن ناحیه محاسبه می‌شود. در این پیاده‌سازی، مقدار میانه به‌عنوان معیار اصلی انتخاب شد، زیرا نسبت به نقاط پرت ناشی از خطای تطبیق استریو، مرزهای تگ و نواحی نامعتبر مقاوم‌تر از میانگین ساده است.

پس از استخراج فاصله و زاویه نسبی هدف، مختصات نسبی هدف در دستگاه مختصات اسکوتر محاسبه می‌شود. سپس برای کاهش نوسان اندازه‌گیری و تخمین سرعت نسبی هدف، از همان ساختار فیلتر کالمن استفاده می‌شود که برای پالایش موقعیت و سرعت موانع به کار رفته است. در این روش، بردار حالت شامل موقعیت طولی، سرعت طولی، موقعیت عرضی و سرعت عرضی هدف نسبت به اسکوتر است و اندازه‌گیری ورودی فیلتر از مختصات دوبعدی هدف تشکیل می‌شود. بنابراین، فیلتر کالمن علاوه بر هموارسازی موقعیت نسبی هدف، مؤلفه‌های سرعت نسبی آن را نیز تخمین می‌زند.

استفاده از فیلتر کالمن برای هدف موجب می‌شود تغییرات لحظه‌ای ناشی از نویز نقشه اختلاف مکانی، جابه‌جایی مرکز تگ در تصویر و ناپایداری تشخیص در چند فریم متوالی کاهش یابد. این موضوع به‌ویژه برای تشکیل مشاهده عامل اهمیت دارد، زیرا عامل تصمیم‌گیر علاوه بر فاصله و زاویه هدف، به مؤلفه‌های سرعت نسبی هدف نیز نیاز دارد. جزئیات مدل حالت، مدل اندازه‌گیری و مراحل پیش‌بینی و تصحیح فیلتر کالمن در پیوست \ref{appendix:kalman_speed} ارائه شده است.

در صورت عدم مشاهده هیچ‌یک از تگ‌های مجاز، خروجی گره تشخیص هدف با پرچم عدم مشاهده منتشر می‌شود. در این حالت، مقدار قابل اعتماد جدیدی برای فاصله، زاویه و سرعت نسبی هدف تولید نمی‌شود. در سطح گره تشخیص، وضعیت هدف به‌صورت نامرئی گزارش می‌شود و مختصات تصویر نیز مقدار نامعتبر می‌گیرد. سپس در مرحله تشکیل مشاهده عامل، از این پرچم برای مشخص کردن وضعیت عدم مشاهده هدف استفاده می‌شود و در صورت نیاز، آخرین موقعیت معتبر مشاهده‌شده برای حفظ پیوستگی محاسبات استفاده می‌شود. جزئیات این رفتار در بخش مربوط به تشکیل مشاهده عامل توضیح داده خواهد شد.

همانظور که در شکل \ref{fig:apriltag_goal_detection_result} هم مشخص است، خروجی نهایی این گره در قالب وضعیت هدف منتشر می‌شود و شامل قابل مشاهده بودن هدف، فاصله تخمین‌زده‌شده، زاویه نسبی هدف، مختصات مرکز تگ در تصویر و مؤلفه‌های سرعت نسبی هدف است. این خروجی یکی از ورودی‌های اصلی زنجیره تشکیل مشاهده عامل محسوب می‌شود، زیرا عامل تصمیم‌گیر برای ناوبری هدف‌محور به فاصله، زاویه و سرعت نسبی هدف نیاز دارد.

\begin{figure}[htbp]
	\centering
	\includegraphics[width=0.85\linewidth]{apriltag_goal_detection_result.png}
	\caption{نمونه تشخیص هدف با استفاده از نشانگر \lr{AprilTag} روی تصویر چپ تصحیح‌شده}
	\label{fig:apriltag_goal_detection_result}
\end{figure}

\subsection{تشخیص موانع و استخراج موقعیت نسبی}
\label{subsec:real_obstacle_detection}

پس از تولید نقشه اختلاف مکانی، مرحله بعدی در لایه ادراک، تشخیص موانع و استخراج موقعیت نسبی آن‌ها نسبت به اسکوتر است. در سامانه عملی، این کار با ترکیب خروجی آشکارساز شیء و نقشه اختلاف مکانی انجام شد. به این ترتیب، آشکارساز، محل تقریبی مانع را در تصویر مشخص می‌کند و نقشه اختلاف مکانی، اطلاعات فاصله‌ای همان ناحیه را فراهم می‌آورد.

روند کلی این پردازش در شکل \ref{fig:real_obstacle_detection_depth_pipeline} نشان داده شده است. مطابق این شکل، تصاویر چپ و راست تصحیح‌شده از یک سو وارد گره تولید نقشه اختلاف مکانی می‌شوند و از سوی دیگر، تصویر چپ تصحیح‌شده برای تشخیص موانع به گره \lr{YOLO} داده می‌شود. سپس خروجی این دو مسیر، یعنی نقشه اختلاف مکانی و کادرهای محدودکننده موانع، در گره تخمین عمق ترکیب می‌شوند تا فاصله، زاویه و موقعیت نسبی هر مانع محاسبه شود.

در این پیاده‌سازی، تشخیص موانع فقط روی تصویر چپ تصحیح‌شده انجام می‌شود. از آنجا که تصاویر ورودی سامانه به‌صورت تک‌رنگ هستند، تصویر خاکستری ابتدا به قالب سه‌کاناله تبدیل می‌شود تا با ورودی مورد انتظار مدل \lr{YOLO} سازگار باشد. سپس تصویر برای کاهش هزینه پردازشی به اندازه ورودی مدل تغییر مقیاس داده می‌شود. پس از اجرای آشکارسازی، مختصات کادرهای محدودکننده دوباره به اندازه اصلی تصویر بازگردانده می‌شوند تا با نقشه اختلاف مکانی هم‌خوانی داشته باشند.

خروجی گره \lr{YOLO} شامل کادر محدودکننده، کلاس مانع، میزان اطمینان تشخیص و شناسه مانع است. در این پژوهش، از این خروجی برای تعیین ناحیه متناظر هر مانع روی نقشه اختلاف مکانی استفاده شد. یعنی \lr{YOLO} وظیفه تشخیص محل مانع در تصویر را بر عهده دارد، اما فاصله مانع مستقیماً از مدل تشخیص به دست نمی‌آید؛ بلکه با نمونه‌برداری از نقشه اختلاف مکانی در محدوده همان کادر محاسبه می‌شود.

برای محاسبه عمق هر مانع، ابتدا کادر محدودکننده آن روی نقشه اختلاف مکانی قرار داده می‌شود. سپس ناحیه داخلی کادر تا حدی کوچک‌تر می‌شود تا اثر پیکسل‌های مرزی، نواحی انسداد و خطاهای تطبیق در اطراف جسم کاهش یابد. پس از آن، پیکسل‌های نامعتبر یا دارای اختلاف مکانی نامناسب حذف می‌شوند. اگر تعداد پیکسل‌های معتبر از حداقل تعیین‌شده بیشتر باشد، یک مقدار نماینده برای اختلاف مکانی مانع محاسبه می‌شود. در حالت پیش‌فرض، از میانه مقادیر معتبر برای گزارش مقدار اختلاف مکانی استفاده شده است.
پس از محاسبه اختلاف مکانی نماینده، فاصله مانع با استفاده از پارامترهای کالیبراسیون دوربین و رابطه عمق-اختلاف مکانی که در فصل چهارم معرفی شد، به دست می‌آید. زاویه نسبی مانع نیز از روی موقعیت افقی مرکز کادر نسبت به مرکز تصویر محاسبه می‌شود. با داشتن فاصله و زاویه، مختصات نسبی مانع در دستگاه مختصات اسکوتر به دست می‌آید. در این دستگاه، مؤلفه طولی نشان‌دهنده فاصله مانع در راستای روبه‌روی اسکوتر و مؤلفه عرضی نشان‌دهنده انحراف جانبی مانع نسبت به محور حرکت است.

خروجی این مرحله برای هر مانع شامل شناسه، کلاس، فاصله پالایش‌شده، فاصله خام، مختصات نسبی، زاویه نسبی و اطلاعات سرعت تخمینی است. این اطلاعات در قالب یک پیام متنی ساختاریافته منتشر می‌شوند و در مراحل بعدی برای تشکیل مشاهده عامل و اعمال منطق ایمنی استفاده می‌شوند. تخمین سرعت نسبی و پالایش زمانی موقعیت موانع در مرحله بعد و با استفاده از فیلتر کالمن انجام می‌شود که جزئیات آن در بخش \ref{subsec:real_obstacle_kalman_filtering} و پیوست \ref{appendix:kalman_speed} ارائه شده است.

\subsection{پالایش و ردیابی موانع با فیلتر کالمن}
\label{subsec:real_obstacle_kalman_filtering}

خروجی خام حاصل از ترکیب آشکارسازی \lr{YOLO} و نقشه اختلاف مکانی، به دلیل نویز تصویر، خطای تطبیق استریو، تغییرات کادر تشخیص و ناپایداری پیکسل‌های معتبر، دارای نوسان است. برای کاهش این نوسان و فراهم‌سازی داده پایدارتر برای عامل تصمیم‌گیر، موقعیت نسبی موانع با استفاده از فیلتر کالمن پالایش شد. مبانی ریاضی فیلتر کالمن، مدل حرکت سرعت ثابت، مدل اندازه‌گیری، پالایش مقادیر پرت عمق و نحوه تخمین سرعت نسبی موانع در پیوست \ref{appendix:kalman_speed} ارائه شده است؛ بنابراین در این بخش تنها نحوه به‌کارگیری آن در سامانه عملی توضیح داده می‌شود.

در پیاده‌سازی، برای هر مانع معتبر یک فیلتر کالمن مستقل تعریف شد. بردار حالت هر مانع شامل موقعیت طولی، سرعت طولی، موقعیت عرضی و سرعت عرضی نسبت به دستگاه مختصات اسکوتر است. اندازه‌گیری ورودی فیلتر نیز از موقعیت دوبعدی مانع تشکیل می‌شود که از ترکیب عمق تخمین‌زده‌شده و زاویه نسبی مانع به دست می‌آید. بنابراین، فیلتر کالمن در این بخش فقط برای هموارسازی عمق استفاده نشده، بلکه برای پالایش موقعیت نسبی دوبعدی و تخمین سرعت نسبی مانع به کار رفته است.

پیش از به‌روزرسانی فیلتر کالمن، عمق خام هر مانع در یک بافر کوتاه‌مدت ذخیره شده و مقادیر پرت آن با استفاده از پالایش مبتنی بر بازه بین‌چارکی \lr{IQR} کاهش داده می‌شود. این پالایش روی مقدار عمق انجام می‌شود، نه روی اختلاف مکانی، زیرا رابطه میان عمق و اختلاف مکانی غیرخطی است. سپس مختصات نسبی مانع با استفاده از عمق پالایش‌شده و زاویه نسبی محاسبه شده و به‌عنوان اندازه‌گیری ورودی فیلتر استفاده می‌شود.

در هر گام زمانی، ابتدا بازه زمانی میان دو مشاهده متوالی محاسبه شده و در ماتریس گذار حالت اعمال می‌شود. سپس مرحله پیش‌بینی و اصلاح فیلتر کالمن انجام می‌گیرد. خروجی نهایی شامل موقعیت پالایش‌شده مانع، عمق پالایش‌شده، زاویه نسبی، مؤلفه‌های سرعت نسبی و همچنین مقدار خام عمق برای مقایسه و عیب‌یابی است. باید توجه داشت که سرعت‌های تخمین‌زده‌شده در این مرحله سرعت نسبی نسبت به اسکوتر هستند و به‌عنوان سرعت مطلق مانع در دستگاه مختصات جهانی تفسیر نمی‌شوند.

پارامترهای فیلتر در پیاده‌سازی عملی با هدف ایجاد تعادل میان نرمی خروجی و پاسخ‌گویی به حرکت موانع تنظیم شدند. در این تنظیمات، نویز فرایند موقعیت و سرعت به‌ترتیب برابر ($0.10$) و ($0.80$)، نویز اندازه‌گیری برابر ($0.10$)، و طول بافر عمق برابر ۸ نمونه در نظر گرفته شد. این مقادیر موجب شدند خروجی نسبت به جهش‌های لحظه‌ای عمق خام مقاوم‌تر شود، بدون آنکه تأخیر زیادی در دنبال‌کردن تغییرات موقعیت مانع ایجاد گردد.


\begin{figure}[htbp]
	\centering
\includegraphics[width=0.95\linewidth]{stereo_vision_pipeline_diagram}
\caption{زنجیره پردازش تصویر استریو، تشخیص موانع و تخمین عمق و موقعیت نسبی موانع در سامانه عملی}
\label{fig:real_obstacle_detection_depth_pipeline}
	
	\vspace{0.4em}
	
	{\small
		\href{https://mohammadabedi1179.github.io/thesis-archive/chapters/ch5/stereo-vision-pipeline-diagram.html}
		{\textbf{مشاهده نسخه تعاملی شکل با کیفیت بالاتر}}
		\qquad
		\qrcode[height=1.2cm]{https://mohammadabedi1179.github.io/thesis-archive/chapters/ch5/stereo-vision-pipeline-diagram.html}
	}
\end{figure}

\subsection{جمع‌بندی بخش بینایی لایه ادراک سامانه عملی}
\label{subsec:real_perception_layer_summary}



\begin{figure}[htbp]
	\centering
\includegraphics[width=0.75\linewidth]{stereo_image_acquisition_pipeline_diagram.png}
\caption{جمع‌بندی زنجیره لایه ادراک در سامانه عملی، از دریافت تصویر استریو تا تولید خروجی‌های مورد نیاز برای برآورد عمق، تشخیص موانع و تشخیص هدف}
\label{fig:real_perception_pipeline_summary}
	
	\vspace{0.4em}
	
	{\small
		\href{https://mohammadabedi1179.github.io/thesis-archive/chapters/ch5/stereo_image_acquisition_pipeline_diagram.html}
		{\textbf{مشاهده نسخه تعاملی شکل با کیفیت بالاتر}}
		\qquad
		\qrcode[height=1.2cm]{https://mohammadabedi1179.github.io/thesis-archive/chapters/ch5/stereo_image_acquisition_pipeline_diagram.html}
	}
\end{figure}

بر اساس توضیحات ارائه‌شده در این بخش، بینایی لایه ادراک سامانه عملی از یک زنجیره چندمرحله‌ای تشکیل شده است که از دریافت تصاویر خام استریو آغاز می‌شود و به تولید خروجی‌های قابل استفاده برای ناوبری ختم می‌گردد. چنان‌که در شکل \ref{fig:real_perception_pipeline_summary} نشان داده شده است، تصاویر دو حسگر دوربین استریو ابتدا در رزبری دریافت می‌شوند. این تصاویر در قالب \lr{YUV420} در دسترس قرار گرفته و با استخراج مؤلفه روشنایی، به تصویر تک‌رنگ \lr{mono8} تبدیل می‌شوند. سپس دو جریان تصویری مستقل از طریق \lr{GStreamer} و بستر شبکه محلی به لپ‌تاپ ارسال می‌گردند. در لپ‌تاپ، تصاویر دریافت‌شده در قالب \lr{ROS2} دریافت شده، تصحیح هندسی روی آن‌ها اعمال می‌شود و تصاویر چپ و راست تصحیح‌شده برای مراحل بعدی در دسترس قرار می‌گیرند.

پس از این مرحله، تصاویر تصحیح‌شده مبنای سه خروجی اصلی لایه ادراک قرار می‌گیرند. خروجی نخست، برآورد اختلاف مکانی و عمق است که برای تخمین فاصله هدف و موانع استفاده می‌شود. خروجی دوم، تشخیص موانع است که با استفاده از آشکارساز دیداری و ترکیب آن با نقشه اختلاف مکانی، موقعیت نسبی و در ادامه سرعت نسبی موانع را فراهم می‌آورد. خروجی سوم، تشخیص هدف است که در سامانه عملی با استفاده از نشانگر \lr{AprilTag} انجام می‌شود. در مجموع، شکل \ref{fig:real_perception_pipeline_summary} نشان می‌دهد که لایه ادراک در سامانه عملی تنها شامل دریافت تصویر نیست، بلکه یک زنجیره یکپارچه از دریافت، تبدیل قالب، انتقال، تصحیح هندسی و استخراج اطلاعات سطح‌بالا را در بر می‌گیرد که در نهایت داده‌های لازم برای تشکیل مشاهده عامل تصمیم‌گیر را فراهم می‌آورد.

\subsection{پردازش داده حسگرهای فراصوت}
\label{subsec:real_ultrasonic_processing}

پنج حسگر \lr{SRF05} در کنار ادراک استریو، موانع نزدیک اطراف اسکوتر را پایش می‌کنند و در نواحی نیازمند واکنش سریع نقش مکمل دارند.

آردوینو برای هر حسگر پالس تحریک ($10~\mu s$) می‌فرستد، زمان پژواک را اندازه می‌گیرد و فاصله را از زمان رفت‌وبرگشت موج صوتی محاسبه می‌کند.

خواندن حسگرها غیرمسدودکننده و نوبتی است تا کنترل موتور و دریافت سریال متوقف نشود. نبود پاسخ معتبر با مقدار (-1) و به‌معنای اندازه‌گیری نامعتبر یا خارج از محدوده گزارش می‌شود.

خروجی، پنج فاصله متناظر با محل نصب است که در مشاهده و ایمنی نزدیک‌برد استفاده می‌شوند. این حسگرها هندسه کامل مانع را نمی‌دهند، اما فاصله چند جهت مشخص را برای افزایش ایمنی فراهم می‌کنند.

\subsection{پردازش داده \lr{IMU}، انکودر و اودومتری نسبی}
\label{subsec:real_imu_encoder_odometry_processing}

برای ثبت و تحلیل حرکت اسکوتر، موقعیت و جهت‌گیری آن به‌صورت نسبی نسبت به نقطه شروع هر آزمون محاسبه می‌شود. این تخمین شامل مختصات صفحه‌ای $x$ و $y$، زاویه جهت‌گیری $\theta$ و طول مسیر طی‌شده است. در پیاده‌سازی انجام‌شده، سرعت انتقالی از داده انکودر چرخ‌ها و جهت‌گیری از واحد اندازه‌گیری اینرسی به دست می‌آید.

استفاده از \lr{IMU} برای تخمین جهت‌گیری به ماهیت رانش لغزشی اسکوتر مربوط است. در مدل سینماتیکی ایده‌آل، نرخ چرخش بدنه از اختلاف سرعت چرخ‌های چپ و راست محاسبه می‌شود؛ بااین‌حال، در اسکوتر چهارچرخ دارای چرخ‌های ثابت، بخشی از دوران چرخ‌ها هنگام پیچیدن صرف لغزش و سایش جانبی می‌شود. در نتیجه، نرخ چرخش حاصل از انکودرها معمولاً با نرخ چرخش واقعی بدنه برابر نیست. به همین دلیل، در حالت عادی زاویه \lr{yaw} حسگر اینرسی منبع اصلی جهت‌گیری است و تخمین انکودری تنها برای مقایسه، عیب‌یابی و حالت پشتیبان استفاده می‌شود.

داده زاویه \lr{yaw} از طریق ارتباط سریال و به‌صورت زاویه اویلر دریافت می‌شود. در آغاز آزمون، چند نمونه متوالی در حالت سکون جمع‌آوری شده و میانگین آن‌ها به‌عنوان مرجع اولیه $\psi_0$ در نظر گرفته می‌شود. ازآنجاکه خروجی خام حسگر در بازه $[-\pi,\pi]$ قرار دارد، اختلاف زاویه نمونه‌های متوالی پیش از انتگرال‌گیری بازگشایی می‌شود تا عبور از مرز $-\pi$ و $+\pi$ موجب پرش ناگهانی نشود. زاویه نسبی مورد استفاده در اودومتری به‌صورت زیر تعریف می‌شود:

\begin{equation}
	\theta_k =
	\operatorname{wrap}
	\left(
	\psi^{\mathrm{unwrap}}_k-\psi_0
	\right)
	\label{eq:real_imu_relative_yaw}
\end{equation}

که در آن $\psi^{\mathrm{unwrap}}_k$ زاویه پیوسته حسگر و تابع
$\operatorname{wrap}(\cdot)$ نگاشت زاویه به بازه $[-\pi,\pi]$ است.

سرعت انتقالی از مقادیر \lr{RPM} چهار چرخ محاسبه می‌شود. به دلیل قرارداد علامت متفاوت موتورهای چپ و راست، ابتدا قدرمطلق \lr{RPM} اندازه‌گیری‌شده در نظر گرفته شده و جهت حرکت هر سمت از علامت آخرین فرمان معتبر همان سمت بازیابی می‌شود. اگر مقادیر \lr{RPM} چرخ‌های جلو و عقب یک سمت به‌ترتیب با
$n_{\mathrm{LF}}$ و $n_{\mathrm{LR}}$ برای سمت چپ و
$n_{\mathrm{RF}}$ و $n_{\mathrm{RR}}$ برای سمت راست نمایش داده شوند، اندازه سرعت دورانی منتخب هر سمت برابر است با:

\begin{equation}
	n_L =
	s_L\min
	\left(
	|n_{\mathrm{LF}}|,
	|n_{\mathrm{LR}}|
	\right),
	\qquad
	n_R =
	s_R\min
	\left(
	|n_{\mathrm{RF}}|,
	|n_{\mathrm{RR}}|
	\right)
	\label{eq:real_odom_side_rpm}
\end{equation}

که در آن $s_L$ و $s_R$ علامت حرکت سمت چپ و راست هستند. استفاده از مقدار کوچک‌تر میان چرخ جلو و عقب، یک انتخاب محافظه‌کارانه برای کاهش اثر چرخ هرزگرد است؛ زیرا در حالت کاهش چسبندگی، چرخ لغزنده معمولاً سریع‌تر از سرعت متناظر با حرکت واقعی زمین می‌چرخد.

با قطر چرخ $D=0.185~m$، محیط چرخ برابر $C=\pi D$ است و سرعت خطی هر سمت از روابط زیر محاسبه می‌شود:

\begin{equation}
	v_L = \frac{n_L C}{60},
	\qquad
	v_R = \frac{n_R C}{60}
	\label{eq:real_odom_side_velocity}
\end{equation}

سرعت خطی خام حاصل از انکودرها نیز برابر میانگین سرعت دو سمت است:

\begin{equation}
	v_{\mathrm{hall}} =
	\frac{v_L+v_R}{2}
	\label{eq:real_odom_hall_velocity}
\end{equation}

علاوه بر انتخاب کمینه سرعت چرخ‌ها، اختلاف سرعت چرخ‌های جلو و عقب هر سمت نیز به‌عنوان شاخص لغزش داخلی همان سمت محاسبه می‌شود. برای سمت چپ، این شاخص به‌صورت زیر است:

\begin{equation}
	\lambda_L =
	\frac{
		\max\left(|n_{\mathrm{LF}}|,|n_{\mathrm{LR}}|\right)
		-
		\min\left(|n_{\mathrm{LF}}|,|n_{\mathrm{LR}}|\right)
	}{
		\max\left(|n_{\mathrm{LF}}|,|n_{\mathrm{LR}}|\right)
	}
	\label{eq:real_odom_front_rear_slip}
\end{equation}

و رابطه مشابهی برای سمت راست استفاده می‌شود. مقدار صفر نشان‌دهنده تطابق سرعت چرخ‌های جلو و عقب و مقادیر بزرگ‌تر بیانگر افزایش اختلاف میان آن‌ها است. این شاخص برای پایش و عیب‌یابی ثبت می‌شود و مستقیماً به‌عنوان ضریب اصلاح سرعت به کار نمی‌رود.

نرخ زاویه‌ای مبتنی بر انکودرها با استفاده از فاصله عرضی چرخ‌ها $b=0.57~m$ محاسبه می‌شود:

\begin{equation}
	\omega_{\mathrm{enc}} =
	\frac{v_R-v_L}{b}
	\label{eq:real_odom_encoder_omega}
\end{equation}

این مقدار در شرایط عادی منبع اصلی جهت‌گیری نیست. هنگامی که داده \lr{IMU} معتبر و تازه باشد، نرخ زاویه‌ای واقعی از تغییر زاویه حسگر در بازه زمانی $\Delta t$ به دست می‌آید:

\begin{equation}
	\omega_{\mathrm{imu},k} =
	\frac{
		\operatorname{wrap}
		\left(
		\theta_k-\theta_{k-1}
		\right)
	}{
		\Delta t
	}
	\label{eq:real_odom_imu_omega}
\end{equation}

در صورت قطع ارتباط \lr{IMU} یا قدیمی‌شدن داده آن، زاویه به‌طور موقت با انتگرال‌گیری از $\omega_{\mathrm{enc}}$ به‌روزرسانی می‌شود:

\begin{equation}
	\theta_k =
	\operatorname{wrap}
	\left(
	\theta_{k-1}
	+
	\omega_{\mathrm{enc},k}\Delta t
	\right)
	\label{eq:real_odom_encoder_fallback}
\end{equation}

برای کاهش بیش‌برآورد سرعت ناشی از سایش جانبی در چرخش‌های رانش لغزشی، یک ضریب اصلاح تجربی از مقایسه نرخ چرخش انکودری و نرخ چرخش اندازه‌گیری‌شده توسط \lr{IMU} محاسبه می‌شود. درصورتی‌که داده \lr{IMU} معتبر باشد، قدرمطلق $\omega_{\mathrm{enc}}$ از آستانه تعیین‌شده بیشتر باشد و علامت دو نرخ زاویه‌ای یکسان باشد، ضریب خام اصلاح به‌صورت زیر محاسبه می‌شود:

\begin{equation}
	k_{\mathrm{raw}} =
	\operatorname{clip}
	\left(
	\frac{\omega_{\mathrm{imu}}}
	{\omega_{\mathrm{enc}}},
	0,1
	\right)
	\label{eq:real_odom_raw_slip_factor}
\end{equation}

در حرکت تقریباً مستقیم، نامعتبر بودن داده \lr{IMU} یا متفاوت بودن علامت دو نرخ زاویه‌ای، ضریب خام برابر یک در نظر گرفته می‌شود؛ زیرا در این شرایط نسبت دو نرخ، شاخص قابل‌اعتمادی برای لغزش نیست. برای جلوگیری از تغییرات ناگهانی، ضریب اصلاح با میانگین‌گیری نمایی هموار می‌شود:

\begin{equation}
	k_k =
	k_{k-1}
	+
	\alpha
	\left(
	k_{\mathrm{raw},k}-k_{k-1}
	\right)
	\label{eq:real_odom_slip_ema}
\end{equation}

که مقدار پیش‌فرض ضریب هموارسازی در پیاده‌سازی برابر $\alpha=0.25$ است. سرعت اصلاح‌شده مورد استفاده در محاسبه موقعیت از رابطه زیر به دست می‌آید:

\begin{equation}
	v_k =
	k_k v_{\mathrm{hall},k}
	\label{eq:real_odom_corrected_velocity}
\end{equation}

این اصلاح بر این فرض تجربی استوار است که نسبتی از دوران چرخ‌ها که به دوران واقعی بدنه تبدیل می‌شود، می‌تواند به‌عنوان تقریب مرتبه اول برای نسبت تبدیل سرعت دورانی چرخ‌ها به جابه‌جایی واقعی نیز استفاده شود. بنابراین، ضریب مذکور یک مدل کامل دینامیکی لغزش نیست و باید به‌عنوان اصلاح تجربی اودومتری تفسیر شود.

برای انتگرال‌گیری موقعیت، زاویه میانی میان دو نمونه متوالی استفاده می‌شود تا خطای تقریب در هنگام چرخش کاهش یابد:

\begin{equation}
	\bar{\theta}*k =
	\frac{\theta*{k-1}+\theta_k}{2}
	\label{eq:real_odom_mid_heading}
\end{equation}

سپس مختصات نسبی اسکوتر مطابق روابط زیر به‌روزرسانی می‌شوند:

\begin{equation}
	x_k =
	x_{k-1}
	+
	v_k\cos\left(\bar{\theta}*k\right)\Delta t,
	\qquad
	y_k =
	y*{k-1}
	+
	v_k\sin\left(\bar{\theta}_k\right)\Delta t
	\label{eq:real_odom_position_update}
\end{equation}

طول مسیر طی‌شده نیز با انتگرال‌گیری از قدرمطلق سرعت اصلاح‌شده محاسبه می‌شود:

\begin{equation}
	d_k =
	d_{k-1}
	+
	|v_k|\Delta t
	\label{eq:real_odom_path_update}
\end{equation}

بازه‌های زمانی نامعتبر یا غیرعادی نیز پیش از انتگرال‌گیری حذف می‌شوند تا توقف یا تأخیر اجرای برنامه موجب جهش موقعیت نشود.

روش مذکور یک اودومتری نسبی و آزمایشگاهی فراهم می‌کند و جایگزین سامانه مکان‌یابی مطلق نیست. دقت آن همچنان به کیفیت داده انکودر، وضعیت تماس چرخ‌ها، ناهمواری سطح، نویز و رانش حسگر اینرسی و خطای تجمعی انتگرال‌گیری وابسته است. همچنین انتخاب سرعت کمتر میان چرخ‌های جلو و عقب و استفاده از نسبت نرخ‌های زاویه‌ای به‌عنوان ضریب اصلاح، تقریب‌های تجربی هستند و ممکن است در شرایطی مانند قفل‌شدن یک چرخ، لغزش نامتقارن شدید یا حرکت روی سطوح ناهمگن معتبر نباشند. بااین‌حال، ترکیب سرعت انتقالی مبتنی بر انکودر با جهت‌گیری مبتنی بر \lr{IMU}، نسبت به اودومتری صرفاً مبتنی بر مدل سینماتیکی چرخ‌ها، تخمین پایدارتر و واقع‌بینانه‌تری از مسیر اسکوتر در چرخش‌های رانش لغزشی فراهم می‌کند.

\subsection{تجمیع داده‌های ادراکی و ساخت مشاهده}
\label{subsec:real_observation_construction}

پس از استخراج خروجی‌های لایه ادراک، داده‌های حاصل باید به قالب مشاهده عامل تبدیل شوند. هدف از این مرحله آن است که مدل آموزش‌دیده در محیط شبیه‌سازی، در سامانه عملی نیز ورودی‌ای با همان ساختار، ترتیب و مقیاس عددی دریافت کند. در این پیاده‌سازی، این کار در گره اجرای مدل انجام می‌شود. این گره داده‌های مربوط به هدف، موانع، حسگرهای فراصوت و اودومتری را از موضوع‌های ارتباطی \lr{ROS2} دریافت کرده و مشاهده نهایی را برای ورودی مدل \lr{LSTM-SAC} می‌سازد.

مشاهده از بردار ثابت ۱۵‌مؤلفه‌ای و توالی موانع تشکیل می‌شود. بردار ثابت وضعیت هدف، حرکت اسکوتر، پنج حسگر فراصوت، فرمان قبلی و وضعیت ایمنی را با ترتیب محیط آموزش نگه می‌دارد:

\begin{equation}
	\begin{aligned}
		\mathbf{o}_{\mathrm{fixed}} =
		[&
		d_g/15,\;
		\beta_g/\pi,\;
		v_{gx}^{\mathrm{rel}}/2,\;
		v_{gy}^{\mathrm{rel}}/2,\;
		v/1,\;
		\omega/1,\;\\
		&
		d_{ufl}/3,\;
		d_{ufr}/3,\;
		d_{uls}/3,\;
		d_{urs}/3,\;
		d_{ub}/3,\;\\
		&
		v_{\mathrm{prev}}/1,\;
		\omega_{\mathrm{prev}}/1,\;
		s_{\mathrm{shield}},\;
		t_{\mathrm{shield}}
		]
	\end{aligned}
	\label{eq:real_fixed_observation_vector}
\end{equation}

در این رابطه، $d_g$ و $\beta_g$ فاصله و زاویه هدف، $v_{gx}^{rel}$ و $v_{gy}^{rel}$ سرعت نسبی هدف، $v$ و $\omega$ سرعت اسکوتر، پنج $d_u$ فاصله‌های فراصوت، $v_{prev}$ و $\omega_{prev}$ فرمان قبلی و دو مؤلفه آخر شدت و مدت فعال‌بودن سپر ایمنی‌اند.

بخش دوم، ماتریس موانع با پنج ویژگی مختصات، فاصله و سرعت نسبی برای هر مانع است:

\begin{equation}
	\mathbf{o}*{obs,i} =
	[
	x_i/10,;
	y_i/10,;
	d_i/10,;
	v*{xi}^{rel}/2,;
	v_{yi}^{rel}/2
	]
	\label{eq:real_obstacle_observation_vector}
\end{equation}

موانع بر اساس فاصله مرتب و نزدیک‌ترین‌ها در ابتدای ماتریس قرار می‌گیرند. تعداد ردیف‌ها باید دقیقاً با بیشینه موانع محیط آموزش یکسان باشد تا شکل ورودی مدل تغییر نکند.

برای جلوگیری از گسست، حافظه‌ای ۱۰ثانیه‌ای از موانع در چارچوب جهانی نگه‌داری می‌شود. موانع محلی با اودومتری به چارچوب جهانی منتقل و هنگام ساخت مشاهده دوباره نسبت به وضعیت فعلی به چارچوب محلی بازگردانده می‌شوند؛ در نتیجه حذف کوتاه‌مدت آشکارسازی تغییر ناگهانی ایجاد نمی‌کند.

بردار ثابت و ماتریس موانع پس از نرمال‌سازی با ضرایب آموزش و محدودشدن به بازه‌های معتبر وارد مدل می‌شوند. حفظ ترتیب و مقیاس ضروری است، زیرا \lr{LSTM-SAC} به قرارداد ورودی حساس است.

\subsection{مدیریت داده‌های نامعتبر}
\label{subsec:real_invalid_data_handling}

در واقعیت، داده حسگری ممکن است ناقص، نویزی یا با تأخیر باشد؛ هدف چند فریم گم شود، مانع موقتاً حذف گردد یا فراصوت نامعتبر برسد. چون عامل با مشاهده ثابت آموزش دیده، این وضعیت‌ها پیش از ورود به مدل مدیریت می‌شوند تا شکل، ترتیب و دامنه ورودی حفظ شود.

فراصوت نامعتبر با (-1) یا داده قدیمی مشخص می‌شود و در هر دو حالت با بیشینه برد معتبر جایگزین می‌گردد تا نبود اندازه‌گیری به‌اشتباه مانع نزدیک تفسیر نشود.

فاصله و زاویه هدف فقط هنگام مشاهده \lr{AprilTag} و عمق معتبر پذیرفته می‌شوند. در نبود هدف، مشاهده تازه‌ای بر پایه آن ساخته نمی‌شود و آخرین مشاهده معتبر برای جلوگیری از جهش نگه داشته می‌شود؛ این راهکار برای گم‌شدن طولانی هدف طراحی نشده است.

تازگی فهرست موانع نیز بررسی می‌شود؛ داده تازه حافظه را اصلاح می‌کند و در نبود آن، وضعیت قبلی تا زمان انقضا می‌ماند. این روش حذف لحظه‌ای و باقی‌ماندن دائمی داده کهنه را هر دو محدود می‌کند.

سرانجام همه مؤلفه‌ها نرمال و محدود می‌شوند. هدف مدیریت داده نامعتبر تغییر سیاست نیست، بلکه حفظ قرارداد ورودی مدل در شرایط واقعی است.

\subsection{اجرای مدل \lr{LSTM-SAC}}
\label{subsec:real_lstm_sac_execution}

مدل آموزش‌دیده روی لپ‌تاپ اجرا می‌شود تا توان کافی برای شبکه عصبی، پردازش تصویر، تشخیص موانع و ارتباط گره‌ها فراهم باشد.

در هر گام، مشاهده وارد سیاست می‌شود و کنش در حالت قطعی تولید می‌گردد تا رفتار آزمون واقعی پایدار و تکرارپذیر باشد. خروجی خام، کنش دوبعدی نرمال‌شده زیر است:

\begin{equation}
	\mathbf{a}*t =
	\begin{bmatrix}
		a_v \
		a*\omega
	\end{bmatrix},
	\qquad
	a_v,a_\omega \in [-1,1]
	\label{eq:real_policy_action}
\end{equation}

در این بردار، $a_v$ و $a_\omega$ به‌ترتیب مؤلفه سرعت خطی و زاویه‌ای‌اند و طبق منطق فصل چهارم مقیاس می‌شوند:

\begin{equation}
	v_{des} = a_v v_{max},
	\qquad
	\omega_{des} = a_\omega \omega_{max}
	\label{eq:real_action_scaling}
\end{equation}

مقادیر ($v_{max}=1.0~m/s$) و ($\omega_{max}=1.0~rad/s$) مطابق محیط آموزش حفظ شدند.

برای کاهش جهش فرمان، خروجی مدل با میانگین‌گیری نمایی از فرمان قبلی و جدید هموار می‌شود:

\begin{equation}
	v_t =
	(1-\alpha)v_{t-1}
	+
	\alpha v_{des}, \quad
	\omega_t =
	(1-\alpha)\omega_{t-1}
	+
	\alpha \omega_{des}
	\label{eq:real_w_and_v_smoothing}
\end{equation}

با $\alpha=0.70$، فرمان همچنان تابع تصمیم جدید است، اما تغییرات لحظه‌ای کاهش می‌یابد و اجرای سخت‌افزاری نرم‌تر می‌شود.

گره، سرعت‌های خطی $v_t$ و زاویه‌ای $\omega_t$ را در پیام \lr{Twist} منتشر می‌کند تا به فرمان چهار موتور تبدیل شوند.

\subsection{تبدیل خروجی مدل به فرمان حرکتی}
\label{subsec:real_action_to_motion_command}

پیام \lr{Twist} شامل $(v,\omega)$ است و برای اسکوتر چهارموتوره باید پیش از اجرا به سرعت هدف هر موتور تبدیل شود.

گره واسط سخت‌افزاری ابتدا سرعت معادل سمت چپ و راست را با مدل رانش تفاضلی محاسبه می‌کند:

\begin{equation}
	v_L = v - \omega \frac{L_{eff}}{2},
	\qquad
	v_R = v + \omega \frac{L_{eff}}{2}
	\label{eq:real_left_and_right_side_velocity}
\end{equation}

در این رابطه، $L_{eff}=0.66108~m$ فاصله مؤثر دو سمت حرکتی است. سپس با قطر چرخ $D=0.185~m$، سرعت خطی هر سمت به \lr{RPM} تبدیل می‌شود:

\begin{equation}
	RPM = \frac{v_{wheel}}{\pi D} \times 60
	\label{eq:real_wheel_speed_to_rpm}
\end{equation}

چرخ جلو و عقب هر سمت هدف یکسان می‌گیرند. به‌دلیل جهت نصب موتورهای راست، علامت فرمان آن‌ها در بردار چهارموتوره معکوس می‌شود:

\begin{equation}
	\mathbf{u}_{rpm} =
	[
	RPM_L,;
	-RPM_R,;
	RPM_L,;
	-RPM_R
	]
	\label{eq:real_four_motor_rpm_command}
\end{equation}

ترتیب مؤلفه‌ها جلو-چپ، جلو-راست، عقب-چپ و عقب-راست است و \lr{RPM} هر موتور برای جلوگیری از فرمان بیش از حد به حداکثر ($150~RPM$) محدود می‌شود.

پس از تشکیل فرمان چهارموتوره، لپ‌تاپ آن را از طریق \lr{UDP} به رزبری پای ارسال می‌کند. رزبری پای فرمان دریافتی را به قالب سریال قابل فهم برای آردوینو تبدیل کرده و از طریق \lr{USB/Serial} به آن منتقل می‌کند. سرانجام، آردوینو با استفاده از بازخورد انکودرها، کنترل سرعت پایین‌سطح هر موتور را انجام داده و سیگنال \lr{PWM} مناسب را برای راه‌اندازهای موتور تولید می‌کند. برای جلوگیری از باقی‌ماندن آخرین فرمان در شرایط خطا، در مسیر ارتباطی فرمان، سازوکار توقف خودکار نیز در نظر گرفته شد. در صورت قطع ارتباط یا دریافت‌نشدن فرمان تازه، فرمان صفر برای موتورهای چهارگانه ارسال می‌شود.
\section{پیاده‌سازی کنترل پایین‌سطح و ایمنی در سامانه عملی}
\label{sec:real_low_level_control_and_safety}

پس از تولید فرمان حرکتی توسط مدل تصمیم‌گیر، لازم است این فرمان به‌صورت قابل اجرا و ایمن به موتورهای اسکوتر اعمال شود. در سامانه عملی، فرمان عامل مستقیماً به راه‌اندازهای موتور ارسال نمی‌شود، بلکه ابتدا به سرعت هدف چهار موتور تبدیل شده و سپس در آردوینو، با استفاده از بازخورد انکودرها، کنترل سرعت پایین‌سطح انجام می‌گیرد. علاوه بر این، برای کاهش خطر برخورد و جلوگیری از ادامه حرکت در شرایط خطا، منطق ایمنی و سازوکار توقف خودکار نیز در مسیر فرمان در نظر گرفته شد.

\subsection{دریافت فرمان و کنترل سرعت موتورهای چهارگانه}
\label{subsec:real_motor_speed_control}

فرمان چهارموتوره از رزبری پای به آردوینو می‌رسد و شامل \lr{RPM} هدف به‌ترتیب جلو-چپ، جلو-راست، عقب-چپ و عقب-راست است. آردوینو هر مقدار را ورودی حلقه کنترل همان موتور قرار می‌دهد.

آردوینو با شمارش پالس انکودر سرعت واقعی هر موتور را تخمین و خطا را از اختلاف هدف و اندازه‌گیری محاسبه می‌کند:

\begin{equation}
	e_i = RPM^{ref}_i - RPM^{meas}_i
	\label{eq:real_motor_rpm_error}
\end{equation}

در این رابطه، $RPM^{ref}_i$، $RPM^{meas}_i$ و $e_i$ به‌ترتیب سرعت هدف، سرعت اندازه‌گیری‌شده و خطای موتور $i$ هستند. خطا، \lr{PWM} و جهت چرخش را تعیین می‌کند تا سرعت واقعی به هدف نزدیک شود؛ بنابراین کنترل حلقه‌بسته است.

\subsection{اعمال منطق ایمنی در مسیر فرمان}
\label{subsec:real_safety_logic_execution}

فرمان عامل پیش از اجرا بررسی می‌شود تا بدون تغییر سیاست یادگرفته‌شده، از نزدیک‌شدن خطرناک یا برخورد جلوگیری گردد. سپر ایمنی همان ساختار فصل چهارم را دارد، اما در عمل مستقیماً سرعت خطی و زاویه‌ای را اصلاح می‌کند.

در این منطق، داده حسگرهای فراصوت برای تشخیص موانع بسیار نزدیک و خروجی تشخیص دیداری برای بررسی موانع روبه‌روی اسکوتر استفاده می‌شود. اگر هیچ مانع نزدیکی در محدوده خطر وجود نداشته باشد، فرمان عامل بدون تغییر اجرا می‌شود. اما در صورت ورود به نواحی احتیاط، خطر یا بحرانی، فرمان نهایی مطابق الگوریتم \ref{alg:real_safety_shield_execution} اصلاح می‌گردد.
\begin{algorithm}[htb]
	\footnotesize
	\caption{اعمال سپر ایمنی فعال روی فرمان حرکتی عامل}
	\label{alg:real_safety_shield_execution}
	\begin{algorithmic}[1]
		\REQUIRE فرمان عامل $\mathbf{u}*{policy}=(v*{policy},\omega_{policy})$، فاصله حسگرهای فراصوت، فاصله نزدیک‌ترین مانع دیداری روبه‌رو
		\ENSURE فرمان نهایی قابل اجرا $\mathbf{u}*{final}=(v*{final},\omega_{final})$
		
		\STATE $v_{final}\gets v_{policy}$ ، $\omega_{final}\gets \omega_{policy}$
		\STATE $d_{front}\gets \min(d_{FL},d_{FR})$
		\STATE $d_{side}\gets \min(d_{LS},d_{RS})$
		\STATE $d_{all}\gets \min(d_{front},d_{side}, d_{back})$
		\STATE $d_{left}\gets \min(d_{FL},d_{LS})$ ، $d_{right}\gets \min(d_{FR},d_{RS})$
		\IF{$d_{all}<0.25~m$ یا $d^{obs}_{front}<0.9~m$}
		\IF{$d_{back}>0.5~m$}
		\STATE $v_{final}\gets -0.15~m/s$
		\ELSE
		\STATE $v_{final}\gets 0$
		\ENDIF
		\IF{$d_{left}\geq d_{right}$}
		\STATE $\omega_{final}\gets +0.8~rad/s$
		\ELSE
		\STATE $\omega_{final}\gets -0.8~rad/s$
		\ENDIF
		\RETURN $\mathbf{u}_{final}$
		\ENDIF
		\IF{$d_{front}<0.50~m$ یا $d^{obs}_{front}<1.5~m$}
		\STATE $v_{final}\gets \min(v_{final},0.10~m/s)$
		\IF{مانع نزدیک‌تر در سمت چپ باشد}
		\STATE $\omega_{final}\gets \omega_{final}-0.5~rad/s$
		\ELSE
		\STATE $\omega_{final}\gets \omega_{final}+0.5~rad/s$
		\ENDIF
		\STATE $\omega_{final}\gets \mathrm{clip}(\omega_{final},-1,1)$
		\RETURN $\mathbf{u}_{final}$
		\ENDIF
		\IF{$d_{front}<0.80~m$ یا $d^{obs}_{front}<2.5~m$ یا $d_{side}<0.50~m$}
		\STATE $v_{final}\gets \min(v_{final},0.40~m/s)$
		\RETURN $\mathbf{u}_{final}$
		\ENDIF
		\RETURN $\mathbf{u}_{policy}$ 
	\end{algorithmic}
	
\end{algorithm}

در سطح احتیاط، سرعت مثبت به $0.40~m/s$ محدود و فرمان چرخشی حفظ می‌شود تا فرصت اصلاح مسیر فراهم باشد. در سطح خطر، حد سرعت به $0.10~m/s$ کاهش و بایاس $0.5~rad/s$ برای چرخش به فضای بازتر افزوده می‌شود.

در سطح بحرانی، اگر عقب آزاد باشد، عقب‌روی $-0.15~m/s$ همراه با چرخش $0.8~rad/s$ به سمت فضای بازتر اعمال می‌شود؛ در غیر این صورت عقب‌روی حذف و چرخش درجا اجرا می‌گردد تا برخورد با مانع پشت سر رخ ندهد.

فرمان قابل اجرا در مجموع به‌صورت زیر انتخاب می‌شود:

\begin{equation}
	\mathbf{u}_{final} =
	\begin{cases}
		\mathbf{u}_{policy}, & \text{وضعیت عادی} \\
		\mathbf{u}_{limited}, & \text{وضعیت احتیاط} \\
		\mathbf{u}_{safe}, & \text{وضعیت خطر یا بحرانی}
	\end{cases}
	\label{eq:real_safe_command_selection}
\end{equation}

در این رابطه، $\mathbf{u}_{policy}$ فرمان عامل، $\mathbf{u}_{limited}$ نسخه محدودشده و $\mathbf{u}_{safe}$ فرمان خروج از خطر است. سپر ایمنی ناظر مسیر فرمان است و فقط هنگام ناسازگاری تصمیم عامل با ایمنی سخت‌افزار، آن را محافظه‌کارانه اصلاح می‌کند.

برای جلوگیری از توقف پایدار در بن‌بست، اگر فرمان نهایی چند گام تقریباً صفر بماند، چرخش اجباری به سمت فضای بازتر و در صورت آزادبودن عقب، عقب‌روی اندکی اعمال می‌شود.

\section{آماده‌سازی و اجرای آزمون‌های عملی}
\label{sec:real_test_preparation_and_execution}

پس از تکمیل پیاده‌سازی و آزمون جداگانه لایه‌های ادراک، تصمیم‌گیری، کنترل و ایمنی، سامانه برای ارزیابی یکپارچه در محیط واقعی آماده شد. این فرایند شامل انتخاب و انتقال مدل نهایی، بررسی سازگاری ورودی‌ها و خروجی‌های سیاست با سامانه عملی، اجرای مرحله‌ای سناریوهای آزمایشی و ثبت هم‌زمان داده‌های موردنیاز بود. در تمام آزمون‌های واقعی، مدل تنها در حالت استنتاج اجرا شد و وزن‌های شبکه ثابت باقی ماندند. در این بخش، مراحل آماده‌سازی و اجرای آزمایش‌ها تشریح می‌شود و تحلیل کمی و کیفی نتایج در فصل ششم ارائه خواهد شد.

\subsection{آماده‌سازی مدل و سامانه پیش از آزمون}
\label{subsec:real_model_system_preparation}

مدل نهایی \lr{LSTM-SAC} پس از تکمیل آموزش و ارزیابی در محیط شبیه‌سازی برای استفاده در سامانه عملی انتخاب شد. فایل ذخیره‌شده مدل همراه با ساختار سیاست سفارشی در گره استنتاج بارگذاری گردید و سازگاری فضای مشاهده، ترتیب مؤلفه‌های ورودی، تعداد بیشینه موانع و حدود فضای عمل با تنظیمات زمان آموزش بررسی شد. این کنترل اهمیت ویژه‌ای داشت؛ زیرا هرگونه تغییر در ابعاد یا ترتیب ورودی‌ها می‌توانست موجب تولید فرمان نامعتبر، حتی در صورت بارگذاری موفق مدل، شود.

پیش از ارسال فرمان به موتورها، عملکرد زنجیره استنتاج به‌صورت غیرفعال بررسی شد. در این حالت، داده‌های واقعی حسگرها دریافت و بردار مشاهده تشکیل می‌شد، اما خروجی سیاست تنها ثبت شده و به لایه کنترل ارسال نمی‌گردید. این آزمون برای بررسی صحت نرمال‌سازی داده‌ها، علامت زاویه هدف، ترتیب موانع، مدیریت داده‌های نامعتبر و قرارگرفتن خروجی‌های خطی و زاویه‌ای در حدود مورد انتظار انجام شد.

در مرحله بعد، ارتباط میان گره استنتاج، پل ارتباطی لپ‌تاپ و رزبری‌پای، میکروکنترلرها و کنترل‌کننده‌های موتور بررسی گردید. ابتدا آزمون جهت چرخش و علامت فرمان‌ها در حالی انجام شد که چرخ‌ها از سطح زمین فاصله داشتند. سپس پاسخ موتورهای چپ و راست به فرمان‌های مستقیم و چرخشی با دامنه محدود ارزیابی شد تا از تطابق قرارداد علامت‌ها، ترتیب موتورها و تبدیل فرمان خطی و زاویه‌ای به سرعت مرجع چرخ‌ها اطمینان حاصل شود.

پس از تأیید عملکرد زنجیره فرمان، آزمون‌های زمینی با سرعت محدود و در فضای کنترل‌شده آغاز شدند. در این مرحله، هموارسازی فرمان، حدود سرعت خطی و زاویه‌ای، کنترل پایین‌سطح موتورها، پایش تازگی داده‌های حسگر و سپر ایمنی فعال بودند. همچنین امکان توقف فوری سامانه توسط اپراتور در تمام مراحل حفظ شد. تنها پس از تأیید عملکرد مستقل هر بخش، اجرای کامل و هم‌زمان زنجیره ادراک، تصمیم‌گیری و کنترل مجاز شد.

\subsection{سناریوهای آزمون و نحوه اجرای آن‌ها}
\label{subsec:real_test_scenarios_execution}

آزمون‌های عملی به‌صورت مرحله‌ای طراحی شدند تا پیش از ارزیابی کامل سامانه، صحت عملکرد هر جزء و ارتباط آن با سایر لایه‌ها بررسی شود. تمام آزمایش‌ها ابتدا با سرعت محدود، فاصله ایمن و در محیط کنترل‌شده انجام شدند. پس از اطمینان از عملکرد حسگرها، پردازش تصویر، کنترل موتورها و سپر ایمنی، زنجیره کامل دریافت داده، تشکیل مشاهده، اجرای سیاست و اعمال فرمان فعال شد.

در سناریوی پایه، هدف مجهز به \lr{AprilTag} در فاصله‌ای مشخص و بدون حضور مانع در مقابل اسکوتر قرار گرفت. هدف این آزمون، بررسی پیوستگی زنجیره ادراک و کنترل، صحت تشخیص هدف، علامت زاویه نسبی، جهت چرخش چرخ‌ها و سازگاری فرمان تولیدشده با موقعیت هدف بود. این سناریو همچنین امکان ارزیابی حرکت تقریباً مستقیم و اصلاح تدریجی خطای زاویه نسبت به هدف را فراهم کرد.

در مرحله بعد، هدف در موقعیتی با زاویه جانبی نسبت به محور طولی اسکوتر قرار گرفت تا توانایی سیاست در تولید هم‌زمان فرمان حرکت و چرخش بررسی شود. در برخی اجراها، موقعیت هدف حین حرکت تغییر داده شد تا واکنش سامانه به تغییر فاصله و زاویه هدف، پیوستگی خروجی سیاست و توانایی آن در دنبال‌کردن هدف متحرک ارزیابی گردد.

در سناریوهای دارای مانع، یک یا چند مانع ایستا یا متحرک در مسیر میان اسکوتر و هدف قرار گرفتند. در این شرایط، آشکارساز دیداری، نقشه اختلاف مکانی و فیلتر کالمن برای تخمین موقعیت، عمق و سرعت نسبی موانع به‌کار گرفته شدند. هدف این آزمون‌ها، ثبت رفتار عامل هنگام دریافت هم‌زمان اطلاعات هدف و مانع و بررسی نحوه ایجاد مصالحه میان پیشروی به سمت هدف و حفظ فاصله ایمن بود.

برای ارزیابی عملکرد سامانه در فواصل نزدیک و نواحی خارج از پوشش مناسب دوربین، موانع به‌صورت جداگانه در محدوده حسگرهای فراصوت جلو، طرفین و عقب قرار گرفتند. در این آزمون‌ها، نحوه محدودسازی یا اصلاح فرمان براساس داده‌های فراصوت و منطق سپر ایمنی بررسی شد. رفتار سامانه در سطوح مختلف هشدار، شامل کاهش سرعت، محدودسازی حرکت در جهت مانع و توقف اضطراری نیز ثبت گردید.

در نهایت، پس از تأیید عملکرد جداگانه اجزا، زنجیره کامل سامانه شامل دریافت تصاویر استریو، تشخیص هدف و موانع، تخمین فاصله و سرعت نسبی، تشکیل ورودی سیاست، اجرای مدل منتخب، هموارسازی و اصلاح ایمن فرمان، تبدیل فرمان به سرعت مرجع چرخ‌ها و کنترل پایین‌سطح موتورها به‌صورت هم‌زمان فعال شد. نتایج حاصل از این آزمون‌های یکپارچه در فصل ششم تحلیل می‌شوند.

\subsection{ثبت داده‌ها برای تحلیل عملکرد}
\label{subsec:real_test_data_logging}

برای تحلیل رفتار سامانه پس از پایان هر آزمون، داده‌های لایه‌های ادراک، تصمیم‌گیری، ایمنی، کنترل و اودومتری به‌صورت هم‌زمان ثبت شدند. این داده‌ها امکان بازسازی روند هر اجرا و مقایسه اطلاعات ورودی سیاست، فرمان تولیدشده و پاسخ واقعی سخت‌افزار را فراهم کردند.

از لایه ادراک، وضعیت مشاهده یا عدم مشاهده \lr{AprilTag}، فاصله و زاویه نسبی هدف، موقعیت و عمق موانع، سرعت نسبی تخمین‌زده‌شده توسط فیلتر کالمن و مقادیر حسگرهای فراصوت ذخیره شدند. این اطلاعات برای بررسی کیفیت تشکیل مشاهده، پایداری تشخیص و شناسایی خطاهای ناشی از فقدان داده یا تخمین نادرست فاصله مورد استفاده قرار گرفتند.

از لایه تصمیم‌گیری و ایمنی، خروجی خام خطی و زاویه‌ای سیاست، فرمان هموارشده، فرمان نهایی پس از اعمال محدودیت‌های ایمنی و وضعیت فعال‌بودن هر یک از سطوح سپر ثبت شدند. این داده‌ها امکان بررسی میزان تغییر فرمان عامل در لایه ایمنی و تعیین زمان و علت مداخله سپر را فراهم کردند.

از لایه کنترل پایین‌سطح، فرمان ارسال‌شده به هر موتور، سرعت مرجع، \lr{RPM} اندازه‌گیری‌شده، خطای کنترل، مقدار اعمال‌شده به محرک و پاسخ واقعی چرخ‌ها ذخیره شدند. همچنین موقعیت نسبی، جهت‌گیری، سرعت خطی خام و اصلاح‌شده، نرخ چرخش، طول مسیر و شاخص‌های لغزش حاصل از ترکیب داده‌های انکودر و \lr{IMU} ثبت گردیدند. این مجموعه داده برای مقایسه فرمان مطلوب سیاست با حرکت واقعی اسکوتر و بررسی تأخیر، عدم تقارن و لغزش چرخ‌ها به‌کار رفت.

در برخی آزمون‌ها، تصاویر تصحیح‌شده دوربین، نتیجه تشخیص هدف، کادرهای آشکارسازی موانع، نقشه اختلاف مکانی و ویدئوی اجرای آزمایش نیز ذخیره شدند. داده‌های عددی و دیداری ثبت‌شده در این مرحله، مبنای ارزیابی کمی و تحلیل کیفی عملکرد سامانه در فصل ششم قرار گرفتند.
\section{جمع‌بندی}
\label{sec:chapter5_summary}

در این فصل، فرایند انتقال چارچوب پیشنهادی از محیط شبیه‌سازی به نمونه واقعی اسکوتر تشریح شد. برای این منظور، ابتدا معماری سامانه عملی در قالب یک ساختار چندلایه شامل ادراک، تصمیم‌گیری، ارتباطات و کنترل پایین‌سطح طراحی گردید و نحوه تبادل داده و فرمان میان رایانه، رزبری‌پای، میکروکنترلرها، حسگرها و عملگرها مشخص شد. سپس اجزای سخت‌افزاری، جانمایی حسگرها، سامانه تغذیه و زیرساخت ارتباطی نمونه واقعی معرفی شدند.

در لایه ادراک، زنجیره پردازش تصاویر استریو، تصحیح تصاویر، تولید نقشه اختلاف مکانی، تشخیص هدف و موانع، تخمین فاصله و سرعت نسبی و پردازش داده حسگرهای فراصوت پیاده‌سازی شد. خروجی این بخش پس از بررسی اعتبار داده‌ها و اعمال نرمال‌سازی لازم، در قالبی سازگار با فضای مشاهده عامل آموزش‌دیده تشکیل گردید. همچنین برای ثبت حرکت اسکوتر، اودومتری نسبی با ترکیب سرعت استخراج‌شده از انکودرهای چرخ و جهت‌گیری اندازه‌گیری‌شده توسط واحد اندازه‌گیری اینرسی محاسبه شد تا اثر خطای مدل سینماتیکی در چرخش‌های رانش لغزشی کاهش یابد.

در لایه تصمیم‌گیری، مدل نهایی \lr{LSTM-SAC} با وزن‌های ثابت بارگذاری و تنها در حالت استنتاج اجرا شد. خروجی خطی و زاویه‌ای سیاست پیش از ارسال به عملگرها از مراحل هموارسازی، محدودسازی و بررسی ایمنی عبور کرد. سپس فرمان حرکتی به سرعت مرجع چرخ‌های چپ و راست تبدیل و از طریق پل‌های ارتباطی به کنترل‌کننده پایین‌سطح موتورها ارسال شد. کنترل‌کننده‌های مستقل سرعت، با استفاده از بازخورد انکودر، وظیفه ردیابی مقادیر مرجع چهار موتور را بر عهده داشتند.

برای کاهش خطر ناشی از خطاهای ادراکی، قطع داده یا نزدیک‌شدن به موانع، سپر ایمنی چندسطحی در مسیر فرمان قرار گرفت. این سازوکار با استفاده هم‌زمان از اطلاعات دوربین، موانع و حسگرهای فراصوت، فرمان عامل را در شرایط خطر محدود یا متوقف می‌کرد. علاوه بر آن، تازگی داده‌ها، حدود خروجی سیاست و امکان توقف فوری توسط اپراتور در تمام مراحل آزمون پایش شد.

در پایان، سامانه به‌صورت مرحله‌ای برای آزمون‌های عملی آماده گردید. ابتدا صحت اجزای منفرد و قرارداد علامت فرمان‌ها بررسی شد و سپس سناریوهای هدف ثابت، هدف جانبی یا متحرک، حضور موانع و شرایط نزدیک‌برد اجرا شدند. داده‌های ادراک، خروجی سیاست، وضعیت سپر ایمنی، فرمان موتورها، پاسخ چرخ‌ها و اودومتری به‌صورت هم‌زمان ثبت شدند. این داده‌ها مبنای تحلیل عملکرد اجزای منفرد و ارزیابی یکپارچه سامانه در فصل ششم قرار می‌گیرند.
